% ============================================================================
% INTRODUCCIÓN: CÓMO USAR ESTE LIBRO
% Va en \frontmatter, por eso el capítulo no lleva número.
% ============================================================================
\chapter{Cómo usar este libro}
\label{cap:introduccion}

% La introducción no tiene número de capítulo: dentro de este grupo las
% secciones, figuras y tablas se numeran 1, 2, 3... (sin el prefijo "0.").
% Al cerrar el grupo (final del archivo) se recupera la numeración normal.
\begingroup
\renewcommand{\thesection}{\arabic{section}}
\renewcommand{\thefigure}{\arabic{figure}}
\renewcommand{\thetable}{\arabic{table}}

\begin{objetivos}[title={\ding{43}\ \ Al terminar esta introducción podrás\ldots}]
\begin{itemize}
  \item Explicar qué es Business Intelligence, qué es Business Analytics y
        cuáles son los cuatro tipos de analítica.
  \item Decidir cuándo conviene usar R y cuándo Excel, Power BI, Tableau o
        Python.
  \item Describir el ciclo de trabajo de un proyecto de BI y saber en qué
        módulo del libro se aprende cada paso.
  \item Leer los bloques de código, las salidas de R y las cajas didácticas
        del libro.
  \item Instalar R, RStudio y los paquetes del curso, y crear tu Proyecto de
        RStudio.
  \item Generar los datos del curso y conocer cada tabla, cada columna y
        cómo se relacionan entre sí.
\end{itemize}
\end{objetivos}

% ----------------------------------------------------------------------------
\section{Bienvenida}
\label{sec:m0-bienvenida}
% ----------------------------------------------------------------------------

Lunes, 9:00 de la mañana. La directora comercial de una cadena de tiendas te
escribe: ``Necesito saber qué sucursal vendió más el último trimestre, qué
productos nos están dejando sin margen y si los descuentos de verdad aumentan
las ventas. Lo presento el jueves''.

Si trabajas con hojas de cálculo ya sabes lo que viene: exportar archivos de
varios sistemas, copiar y pegar columnas, corregir fechas que se leyeron mal,
armar tablas dinámicas, hacer gráficas una por una\ldots{} y el lunes
siguiente repetir todo desde cero porque llegaron datos nuevos. Este libro te
enseña a hacer ese mismo trabajo con \textbf{R} y \textbf{RStudio}: escribes
una sola vez las instrucciones en un archivo de código (un \termino{script})
y, a partir de entonces, el análisis completo se repite en segundos, sin
errores de copiar y pegar y con gráficas de calidad profesional.

Es un \textbf{curso práctico}. Todo gira alrededor de una cadena minorista
(\termino{retail}) ficticia con diez sucursales en México: sus ventas de
2023, sus clientes, su catálogo de productos, sus empleados y sus tiendas.
Con esos datos resolverás, paso a paso, las mismas preguntas que se hacen a
diario en las áreas comerciales, de finanzas, de operaciones y de marketing
de cualquier empresa. Cada ejemplo sigue el mismo patrón: una pregunta de
negocio, el código comentado línea por línea, la salida real de R y la
interpretación (``¿qué nos dice esto?''). Después te toca a ti.

\subsection{Lo que sabrás hacer al terminar}

Al final del libro podrás hacer el trabajo diario de un analista de
Business Intelligence:

\begin{itemize}
  \item \textbf{Reunir datos} de archivos CSV, libros de Excel y bases de
        datos SQL, y \textbf{unir} tablas de ventas, clientes, productos y
        tiendas.
  \item \textbf{Limpiar y transformar} datos con scripts reproducibles: lo
        que hoy te toma una hora de trabajo manual se ejecutará con un clic.
  \item \textbf{Calcular indicadores} (ventas, ticket promedio, margen,
        crecimiento, participación) por tienda, mes, categoría o vendedor.
  \item \textbf{Crear gráficas} claras y con estilo corporativo, listas para
        una presentación ejecutiva.
  \item \textbf{Explorar datos con estadística}: detectar valores atípicos,
        medir relaciones y comprobar si una diferencia es real o casualidad.
  \item \textbf{Construir dashboards interactivos} con Shiny que tus
        compañeros pueden usar desde el navegador.
  \item \textbf{Segmentar clientes y hacer pronósticos} con técnicas de
        machine learning.
  \item \textbf{Automatizar reportes} en HTML, Word y PDF que se actualizan
        solos cuando llegan datos nuevos.
  \item \textbf{Integrarlo todo} en un proyecto final: un dashboard
        ejecutivo de \textit{retail analytics}.
\end{itemize}

\subsection{¿Para quién es este libro y qué necesitas?}

No necesitas saber programar. Basta con que manejes una computadora y tengas
nociones de hoja de cálculo (qué es una fila, una columna, una suma). Si ya
programas en otro lenguaje, avanzarás rápido por el primer módulo; si ya
usas R, aprovecharás sobre todo la parte de negocio: indicadores,
dashboards, reportes automáticos y modelos aplicados.

Necesitas una computadora con Windows 10/11, macOS o Linux, al menos 4\,GB
de memoria RAM (8\,GB es más cómodo), unos 2\,GB libres en disco, conexión a
internet para instalar el software y, sobre todo, constancia: el plan
sugerido de la sección~\ref{sec:m0-plan} pide unas 8 horas por semana
durante 10 semanas. Todo el software que usaremos es gratuito.

% ----------------------------------------------------------------------------
\section{¿Qué es Business Intelligence?}
\label{sec:m0-que-es-bi}
% ----------------------------------------------------------------------------

\begin{concepto}{Business Intelligence y Business Analytics}
\termino{Business Intelligence} (BI, inteligencia de negocios) es el
conjunto de procesos, herramientas y prácticas que convierten los datos que
genera una empresa ---ventas, inventarios, clientes, nómina--- en
información confiable para tomar decisiones. Sus productos típicos son
reportes, indicadores y tableros que responden \emph{qué pasó} y \emph{cómo
vamos}.

\termino{Business Analytics} (BA, analítica de negocios) va un paso más
allá: usa estadística y modelos para explicar \emph{por qué} pasó algo,
anticipar \emph{qué pasará} y recomendar \emph{qué hacer}.

En la práctica la frontera es difusa y muchas empresas llaman ``BI'' a
todo. En este libro aprenderás las dos cosas.
\end{concepto}

Dos términos que usarás en todo el libro. Un \termino{KPI} (\textit{Key
Performance Indicator}, indicador clave de desempeño) es una medida numérica
que dice si el negocio avanza hacia un objetivo: ventas del mes, ticket
promedio, margen bruto, porcentaje de clientes que regresan. Un
\termino{dashboard} (tablero) es una pantalla que reúne varios KPI y
gráficas para dar un vistazo rápido al estado del negocio.

\subsection{Los cuatro tipos de analítica}

Las preguntas que se hacen con datos se suelen agrupar en cuatro tipos de
analítica, de menor a mayor complejidad y, en general, de menor a mayor
valor para la decisión (figura~\ref{fig:m0-analitica}). No es una escalera
que se sube una sola vez: una empresa madura usa los cuatro tipos al mismo
tiempo, y la analítica descriptiva, bien hecha, sigue siendo la base de
todo.

La tabla~\ref{tab:m0-analitica} aterriza cada tipo con ejemplos de nuestra
cadena de tiendas y te dice en qué módulo lo aprenderás.

\begin{table}[H]
\centering\small
\caption{Los cuatro tipos de analítica aplicados a una cadena de tiendas.}
\label{tab:m0-analitica}
\begin{tabularx}{\textwidth}{@{}>{\raggedright\arraybackslash}p{2.1cm}
    >{\raggedright\arraybackslash}X
    >{\raggedright\arraybackslash}p{3.7cm}
    >{\raggedright\arraybackslash}p{1.5cm}@{}}
\toprule
\textbf{Tipo} & \textbf{Ejemplo en nuestra cadena} &
\textbf{Técnicas en R} & \textbf{Módulo} \\
\midrule
Descriptiva &
  Ventas totales por mes y por tienda en 2023; los 10 productos más
  vendidos; ticket promedio por método de pago. &
  Resúmenes con \paquete{dplyr}, gráficas con \paquete{ggplot2}. &
  \ref{cap:manipulacion}, \ref{cap:visualizacion} \\
\addlinespace
Diagnóstica &
  ¿Bajaron las ventas de una sucursal porque hubo menos compras o porque
  cada compra fue más pequeña? ¿Los clientes premium gastan más? &
  Comparación de grupos, correlación, pruebas de hipótesis, regresión. &
  \ref{cap:eda} \\
\addlinespace
Predictiva &
  Pronóstico de ventas del próximo trimestre; probabilidad de que un
  cliente deje de comprar. &
  Series de tiempo (\paquete{forecast}), árboles y bosques aleatorios. &
  \ref{cap:ml} \\
\addlinespace
Prescriptiva &
  Qué productos reabastecer primero; a qué segmento de clientes enviar la
  promoción. &
  Reglas de decisión sobre predicciones y segmentos, comparación de
  escenarios. &
  \ref{cap:ml}, \ref{cap:proyecto} \\
\bottomrule
\end{tabularx}
\end{table}

\begin{figure}[H]
\centering
\begin{tikzpicture}[
  escalon/.style={rounded corners=3pt, minimum width=3.45cm,
    minimum height=1.25cm, align=center, text=white,
    font=\sffamily\small, anchor=south west},
  nota/.style={text width=3.3cm, align=center, font=\sffamily\scriptsize,
    text=gristexto, anchor=north},
]
  % Ejes
  \draw[-{Stealth[length=3mm]}, gristexto, thick] (0,0) -- (15.4,0)
    node[below left, font=\sffamily\small, text=gristexto]
    {Complejidad};
  \draw[-{Stealth[length=3mm]}, gristexto, thick] (0,0) -- (0,6.4)
    node[above right, font=\sffamily\small, text=gristexto]
    {Valor para la decisión};
  % Escalones
  \node[escalon, fill=secundario] (d1) at (0.35,0.35)
    {\textbf{Descriptiva}\\¿Qué pasó?};
  \node[escalon, fill=primario] (d2) at (4.1,1.75)
    {\textbf{Diagnóstica}\\¿Por qué pasó?};
  \node[escalon, fill=morado] (d3) at (7.85,3.15)
    {\textbf{Predictiva}\\¿Qué pasará?};
  \node[escalon, fill=acento] (d4) at (11.6,4.55)
    {\textbf{Prescriptiva}\\¿Qué debemos hacer?};
  % Notas bajo cada escalón
  \node[nota] at ([yshift=-2pt]d2.south -| d2.center)
    {Comparaciones, correlación, pruebas de hipótesis};
  \node[nota] at ([yshift=-2pt]d3.south -| d3.center)
    {Pronósticos, modelos de clasificación};
  \node[nota] at ([yshift=-2pt]d4.south -| d4.center)
    {Reglas de decisión, escenarios, optimización};
  \node[nota, anchor=south] at ([yshift=2pt]d1.north)
    {Reportes, KPI, tableros};
\end{tikzpicture}
\caption{Los cuatro tipos de analítica. Cada escalón responde una pregunta
más difícil y, por lo general, más valiosa para la decisión.}
\label{fig:m0-analitica}
\end{figure}

Siendo honestos: la analítica prescriptiva ``completa'' (optimización
matemática de precios, rutas o inventarios) es un campo propio que va más
allá de este libro. Aquí harás la versión que más se usa en el día a día:
convertir segmentos y pronósticos en reglas de decisión claras y comparar
escenarios.

\subsection{Roles de datos en una empresa}

Alrededor de los datos trabajan varios perfiles. Conocerlos te ayuda a
entender qué parte del trabajo te toca y con quién colaboras
(tabla~\ref{tab:m0-roles}).

\begin{table}[H]
\centering\small
\caption{Roles de datos más comunes y la pregunta típica de cada uno.}
\label{tab:m0-roles}
\begin{tabularx}{\textwidth}{@{}>{\raggedright\arraybackslash}p{2.3cm}
    >{\raggedright\arraybackslash}X
    >{\raggedright\arraybackslash}p{3.1cm}
    >{\raggedright\arraybackslash}p{3.6cm}@{}}
\toprule
\textbf{Rol} & \textbf{Qué hace} & \textbf{Herramientas} &
\textbf{Pregunta típica} \\
\midrule
Analista de BI &
  Traduce preguntas de negocio en indicadores; construye reportes y
  dashboards periódicos. &
  SQL, Excel, Power BI o Tableau, R. &
  ¿Cómo vamos contra la meta del mes? \\
\addlinespace
Analista de datos &
  Responde preguntas puntuales con análisis a la medida y estadística. &
  SQL, R o Python, Excel. &
  ¿Qué tipo de cliente responde mejor a los descuentos? \\
\addlinespace
Científico de datos &
  Construye modelos predictivos y diseña experimentos. &
  R o Python, machine learning. &
  ¿Qué clientes tienen alta probabilidad de irse? \\
\addlinespace
Ingeniero de datos &
  Construye y mantiene las ``tuberías'' que llevan los datos de los
  sistemas al almacén de datos. &
  SQL, Python, herramientas de ETL, servicios en la nube. &
  ¿Llegaron completos y a tiempo los datos de ventas de ayer? \\
\bottomrule
\end{tabularx}
\end{table}

\begin{negocio}{una persona, varios sombreros}
En una empresa grande cada rol es un equipo distinto. En una empresa mediana
o pequeña es común que una sola persona extraiga los datos, los limpie, haga
el análisis y presente el dashboard. Este libro te forma principalmente como
\textbf{analista de BI y de datos}, con una introducción sólida a la ciencia
de datos (Módulo~\ref{cap:ml}) y a la parte de ingeniería que todo analista
necesita: bases de datos, SQL y procesos ETL (Módulo~\ref{cap:bases-datos}).
\end{negocio}

\subsection{Preguntas de negocio que resuelve BI}

BI no empieza con los datos, sino con una pregunta. La
tabla~\ref{tab:m0-preguntas} reúne preguntas típicas por área de la empresa y
la tabla del curso con la que podrás responderlas.

\begin{table}[H]
\centering\small
\caption{Preguntas de negocio por área y datos del curso que las responden.}
\label{tab:m0-preguntas}
\begin{tabularx}{\textwidth}{@{}>{\raggedright\arraybackslash}p{2.4cm}
    >{\raggedright\arraybackslash}X
    >{\raggedright\arraybackslash}p{3.3cm}@{}}
\toprule
\textbf{Área} & \textbf{Preguntas típicas} & \textbf{Datos del curso} \\
\midrule
Ventas &
  ¿Qué tiendas venden más? ¿Cómo evolucionan las ventas mes a mes? ¿Qué
  día de la semana vendemos más? &
  \archivo{ventas\_retail}, \archivo{transacciones} \\
\addlinespace
Clientes &
  ¿Quiénes son los mejores clientes? ¿Los clientes premium compran más?
  ¿Qué segmentos existen? &
  \archivo{clientes} + ventas \\
\addlinespace
Productos e inventario &
  ¿Qué productos tienen margen negativo? ¿Cuáles están por debajo del stock
  mínimo? ¿Qué categoría aporta más? &
  \archivo{productos} + ventas \\
\addlinespace
Tiendas &
  ¿Cuánto vende cada tienda por metro cuadrado? ¿Las tiendas más antiguas
  venden más? &
  \archivo{tiendas} + ventas \\
\addlinespace
Recursos humanos &
  ¿Cómo se distribuyen los salarios por departamento, puesto y nivel
  educativo? ¿Qué vendedor vende más? &
  \archivo{empleados} + ventas \\
\addlinespace
Finanzas &
  ¿Cuánto dinero dejamos de cobrar por descuentos? ¿Cómo cambia el ingreso
  por trimestre? &
  \archivo{ventas\_retail} \\
\addlinespace
Marketing &
  ¿Qué método de pago prefieren los clientes? ¿A qué segmento dirigimos la
  próxima campaña? &
  \archivo{transacciones}, \archivo{clientes} \\
\bottomrule
\end{tabularx}
\end{table}

\begin{hazlotu}
Escribe tres preguntas que se hagan en tu trabajo (o en un negocio que
conozcas) y clasifícalas como descriptiva, diagnóstica, predictiva o
prescriptiva. Guárdalas: al terminar el libro podrás responderlas con R.
\end{hazlotu}

% ----------------------------------------------------------------------------
\section{¿Por qué R y RStudio para BI?}
\label{sec:m0-por-que-r}
% ----------------------------------------------------------------------------

\termino{R} es un lenguaje de programación gratuito y de código abierto
creado específicamente para el análisis de datos y la estadística (nació en
los años noventa en la Universidad de Auckland, Nueva Zelanda). Su fuerza
está en sus \termino{paquetes}: extensiones gratuitas que agregan funciones
nuevas. En el repositorio oficial, \termino{CRAN}, hay más de 20\,000. Entre
ellos destaca el \termino{tidyverse}, una colección de paquetes con una
forma de escribir código muy legible, casi como una receta: ``toma las
ventas, filtra 2023, agrupa por tienda, suma el total''.

\begin{concepto}{R y RStudio no son lo mismo}
\textbf{R} es el motor: el programa que ejecuta los cálculos.
\textbf{RStudio} es el tablero del coche: un programa (un \termino{IDE},
entorno de desarrollo integrado) que te da un editor de código, una
consola, un visor de gráficas, un explorador de archivos y la ayuda, todo en
una ventana. RStudio lo desarrolla la empresa Posit y necesita que R esté
instalado. Tú trabajarás siempre dentro de RStudio.
\end{concepto}

\subsection{Comparación honesta con otras herramientas}

R no es la única herramienta de BI ni la mejor para todo. La
tabla~\ref{tab:m0-herramientas} resume fortalezas y limitaciones de las
opciones más comunes.

\begin{table}[H]
\centering\small
\caption{R frente a otras herramientas de análisis de datos.}
\label{tab:m0-herramientas}
\begin{tabularx}{\textwidth}{@{}>{\raggedright\arraybackslash}p{2.25cm}
    >{\raggedright\arraybackslash}X
    >{\raggedright\arraybackslash}X
    >{\raggedright\arraybackslash}p{3.3cm}@{}}
\toprule
\textbf{Herramienta} & \textbf{Fortalezas} & \textbf{Limitaciones} &
\textbf{Úsala cuando\ldots} \\
\midrule
Excel &
  Está en todas las empresas; todo mundo lo sabe leer; ideal para
  cálculos rápidos y tablas dinámicas. &
  Máximo 1\,048\,576 filas por hoja; los procesos de copiar y pegar son
  difíciles de repetir y de auditar; errores de fórmula silenciosos;
  estadística limitada. &
  El análisis es pequeño y único, o necesitas entregar resultados a
  usuarios de negocio. \\
\addlinespace
Power BI y Tableau &
  Dashboards interactivos arrastrando y soltando; publicación con
  permisos para toda la organización; muchos conectores; actualización
  programada. &
  Licencias de pago para compartir; transformaciones complejas y
  estadística avanzada son difíciles; lenguajes propios (DAX, cálculos)
  con su propia curva. &
  La empresa ya lo tiene y necesitas un tablero corporativo de consulta
  diaria para muchas personas. \\
\addlinespace
Python &
  Lenguaje de propósito general; líder en machine learning a gran escala,
  deep learning e integración con otros sistemas. &
  Para estadística y gráficas de reporte suele requerir más código y más
  librerías distintas; configurar el entorno es más confuso al principio. &
  El modelo irá a producción dentro de una aplicación o tu equipo ya
  trabaja en Python. \\
\addlinespace
R y RStudio &
  Diseñado para analizar datos; código legible con el tidyverse; gráficas
  de calidad con \paquete{ggplot2}; reportes reproducibles; dashboards con
  Shiny; gratuito. &
  Hay que programar (curva inicial); los datos se cargan en la memoria
  RAM; se usa menos para desarrollar software; publicar dashboards
  Shiny requiere un servidor. &
  Necesitas limpiar, analizar, modelar y reportar de forma reproducible y
  automatizable. \\
\bottomrule
\end{tabularx}
\end{table}

La gran ventaja de R para BI se resume en una palabra:
\termino{reproducible}. Un análisis es reproducible cuando cualquier
persona ---incluido tú dentro de seis meses--- puede volver a ejecutarlo y
obtener exactamente el mismo resultado. En Excel, los pasos viven en tu
memoria (``primero filtré, luego copié esta columna\ldots''); en R, viven en
un script que se puede revisar, corregir, compartir y volver a ejecutar con
los datos del mes siguiente.

\begin{consejo}
R \textbf{complementa} a tus otras herramientas; no las sustituye. Un
analista eficaz usa cada una para lo que hace mejor. Con R puedes exportar
resultados a Excel para el área de finanzas (Módulo~\ref{cap:fundamentos}),
alimentar un tablero de Power BI (que puede ejecutar scripts de R como
fuente de datos) o generar el reporte en Word que pide dirección
(Módulo~\ref{cap:reportes}). Aprender R no te obliga a abandonar nada.
\end{consejo}

% ----------------------------------------------------------------------------
\section{El ciclo de trabajo de un proyecto de BI}
\label{sec:m0-ciclo}
% ----------------------------------------------------------------------------

Todo proyecto de BI, desde un reporte semanal hasta un modelo de
pronóstico, sigue un ciclo parecido (figura~\ref{fig:m0-ciclo}). Es un
ciclo porque la medición de resultados casi siempre genera una pregunta
nueva. El libro está organizado para que aprendas cada paso; debajo de cada
uno verás el módulo donde lo practicas.

\begin{figure}[H]
\centering
\begin{tikzpicture}[
  paso/.style={draw=#1, fill=#1!7, line width=0.9pt, rounded corners=4pt,
    text width=3.25cm, align=center, inner sep=4pt,
    font=\sffamily\small},
  flecha/.style={-{Stealth[length=2.6mm]}, line width=1pt,
    gristexto!70},
]
  \node[paso=acento] (p1) at (90:6.0cm and 3.9cm)
    {\textbf{1. Pregunta de negocio}\\[1pt]
     {\scriptsize\color{gristexto} ¿Qué hay que decidir?}\\
     {\scriptsize\color{gristexto} Introducción y Módulo \ref{cap:proyecto}}};
  \node[paso=secundario] (p2) at (38.57:6.0cm and 3.9cm)
    {\textbf{2. Obtener los datos}\\[1pt]
     {\scriptsize\color{gristexto} CSV, Excel, SQL}\\
     {\scriptsize\color{gristexto} Módulos \ref{cap:fundamentos} y \ref{cap:bases-datos}}};
  \node[paso=primario] (p3) at (-12.86:6.0cm and 3.9cm)
    {\textbf{3. Limpiar y transformar (ETL)}\\[1pt]
     {\scriptsize\color{gristexto} Filtrar, unir, resumir}\\
     {\scriptsize\color{gristexto} Módulos \ref{cap:manipulacion} y \ref{cap:bases-datos}}};
  \node[paso=morado] (p4) at (-64.29:6.0cm and 3.9cm)
    {\textbf{4. Analizar}\\[1pt]
     {\scriptsize\color{gristexto} Estadística y modelos}\\
     {\scriptsize\color{gristexto} Módulos \ref{cap:eda} y \ref{cap:ml}}};
  \node[paso=exito] (p5) at (-115.71:6.0cm and 3.9cm)
    {\textbf{5. Visualizar y comunicar}\\[1pt]
     {\scriptsize\color{gristexto} Gráficas, dashboards, reportes}\\
     {\scriptsize\color{gristexto} Módulos \ref{cap:visualizacion}, \ref{cap:shiny} y
      \ref{cap:reportes}}};
  \node[paso=acento] (p6) at (-167.14:6.0cm and 3.9cm)
    {\textbf{6. Tomar la decisión}\\[1pt]
     {\scriptsize\color{gristexto} Junto con el negocio}\\
     {\scriptsize\color{gristexto} Módulo \ref{cap:proyecto}}};
  \node[paso=secundario] (p7) at (141.43:6.0cm and 3.9cm)
    {\textbf{7. Medir el resultado}\\[1pt]
     {\scriptsize\color{gristexto} KPI antes y después}\\
     {\scriptsize\color{gristexto} Módulos \ref{cap:shiny} y \ref{cap:reportes}}};
  % Flechas del ciclo
  \foreach \a/\b in {p1/p2, p2/p3, p3/p4, p4/p5, p5/p6, p6/p7, p7/p1}
    \draw[flecha] (\a) to[bend left=12] (\b);
  % Texto central
  \node[align=center, font=\sffamily, text=primario] at (0,0.1)
    {\textbf{Ciclo de un}\\\textbf{proyecto de BI}\\[3pt]
     {\scriptsize\color{gristexto}cada resultado genera}\\
     {\scriptsize\color{gristexto}una pregunta nueva}};
\end{tikzpicture}
\caption{El ciclo de trabajo de un proyecto de BI y el módulo del libro
donde se aprende cada paso.}
\label{fig:m0-ciclo}
\end{figure}

\begin{ejemplo}{el ciclo completo en una pregunta real}
La directora comercial pregunta: \textit{``¿Los descuentos aumentan las
unidades vendidas o solo nos quitan margen?''}. Así se vería el ciclo:
\begin{enumerate}
  \item \textbf{Pregunta:} ¿las ventas con 15 o 20\,\% de descuento venden
        más unidades que las ventas sin descuento?
  \item \textbf{Datos:} \archivo{ventas\_retail.csv}, columnas
        \codigo{descuento\_pct}, \codigo{cantidad} y \codigo{total}.
  \item \textbf{Limpiar:} comprobar que \codigo{total} sea igual a
        \codigo{subtotal} menos \codigo{descuento\_monto} y que no haya
        valores faltantes.
  \item \textbf{Analizar:} comparar las unidades promedio por nivel de
        descuento y probar si la diferencia es estadísticamente
        significativa (Módulo~\ref{cap:eda}).
  \item \textbf{Comunicar:} una gráfica de barras con las unidades
        promedio por nivel de descuento y una frase con la conclusión.
  \item \textbf{Decidir:} si los descuentos no mueven las unidades, se
        reducen a las fechas clave.
  \item \textbf{Medir:} comparar unidades y margen del trimestre siguiente
        contra el anterior.
\end{enumerate}
\end{ejemplo}

\begin{hazlotu}
Piensa en una decisión reciente de tu trabajo (un precio, una compra, una
contratación). Escribe cómo habría sido cada uno de los siete pasos del
ciclo. ¿Qué datos habrías necesitado? ¿Cómo habrías medido el resultado?
\end{hazlotu}

% ----------------------------------------------------------------------------
\section{Cómo está organizado el libro}
\label{sec:m0-organizacion}
% ----------------------------------------------------------------------------

El libro tiene cuatro partes que van de principiante a avanzado, más dos
apéndices (tabla~\ref{tab:m0-modulos}). Cada módulo se apoya en los
anteriores, así que te recomendamos seguirlos en orden, al menos la primera
vez.

\begin{table}[H]
\centering\small
\caption{Partes, módulos, nivel y horas estimadas de estudio y práctica.}
\label{tab:m0-modulos}
\begin{tabularx}{\textwidth}{@{}>{\raggedright\arraybackslash}p{2.3cm}
    c >{\raggedright\arraybackslash}X
    >{\raggedright\arraybackslash}p{2.0cm} r@{}}
\toprule
\textbf{Parte} & \textbf{Módulo} & \textbf{Tema} & \textbf{Nivel} &
\textbf{Horas} \\
\midrule
\multirow{3}{2.3cm}{\raggedright I.~Nivel principiante} &
  \ref{cap:fundamentos} & Fundamentos de R y RStudio: objetos, vectores,
  tablas, funciones, importar y exportar CSV y Excel. & Principiante & 8 \\
& \ref{cap:manipulacion} & Manipulación de datos con \paquete{dplyr} y
  \paquete{tidyr}: filtrar, crear columnas, agrupar, resumir, unir y
  pivotar. & Principiante & 10 \\
& \ref{cap:visualizacion} & Visualización con \paquete{ggplot2}: barras,
  líneas, dispersión, distribuciones y tema corporativo. & Principiante &
  8 \\
\midrule
\multirow{3}{2.3cm}{\raggedright II.~Nivel intermedio} &
  \ref{cap:eda} & Análisis exploratorio y estadística: resúmenes,
  distribuciones, atípicos, correlación, pruebas y regresión. &
  Intermedio & 8 \\
& \ref{cap:bases-datos} & Bases de datos y SQL: SQLite desde R, consultas
  SQL, \paquete{dbplyr} y procesos ETL. & Intermedio & 6 \\
& \ref{cap:shiny} & Dashboards interactivos con Shiny y
  \paquete{shinydashboard}. & Intermedio & 8 \\
\midrule
\multirow{2}{2.3cm}{\raggedright III.~Nivel avanzado} &
  \ref{cap:ml} & Machine learning para BI: segmentación, clasificación,
  regresión y series de tiempo. & Avanzado & 10 \\
& \ref{cap:reportes} & Reportes automáticos con R Markdown en HTML, Word y
  PDF. & Avanzado & 6 \\
\midrule
IV.~Proyecto integrador & \ref{cap:proyecto} & Proyecto final: dashboard
  ejecutivo de \textit{retail analytics}, de los datos crudos al reporte.
  & Integrador & 12 \\
\midrule
Apéndices & \ref{ap:soluciones} & Soluciones completas de todos los
  ejercicios. & --- & --- \\
& \ref{ap:referencia} & Referencia rápida de funciones y atajos. & --- &
  --- \\
\midrule
\multicolumn{4}{@{}l}{\textbf{Total aproximado}} & \textbf{76} \\
\bottomrule
\end{tabularx}
\end{table}

\subsection{Estructura de cada módulo}

Todos los módulos siguen la misma estructura, para que sepas siempre dónde
estás:

\begin{enumerate}
  \item \textbf{Objetivos} y un escenario de negocio que motiva el tema.
  \item \textbf{Secciones de contenido}: cada idea se presenta con un
        concepto breve, un ejemplo con código comentado, la salida real de
        R, su interpretación y un reto ``Hazlo tú''.
  \item \textbf{Casos prácticos} completos, de la pregunta a la conclusión.
  \item \textbf{Errores comunes}: los mensajes de error que verás con más
        frecuencia y cómo resolverlos.
  \item \textbf{Resumen} con lo aprendido y una tabla de funciones clave.
  \item \textbf{Ejercicios} graduados por nivel (Básico, Intermedio,
        Avanzado y Reto). Cada ejercicio indica la página de su solución
        en el Apéndice~\ref{ap:soluciones}.
\end{enumerate}

\subsection{Plan de estudio sugerido de 10 semanas}
\label{sec:m0-plan}

Si estudias por tu cuenta, el plan de la tabla~\ref{tab:m0-plan} reparte
unas 80 horas (los módulos más esta introducción y la planeación del
proyecto) en 10 semanas de 6 a 10 horas. Cada semana termina con un
\textbf{entregable}: algo concreto que puedes guardar en tu portafolio.

\begin{table}[H]
\centering\small
\caption{Plan de estudio sugerido de 10 semanas.}
\label{tab:m0-plan}
\begin{tabularx}{\textwidth}{@{}c
    >{\raggedright\arraybackslash}p{5.2cm}
    >{\raggedright\arraybackslash}X r@{}}
\toprule
\textbf{Semana} & \textbf{Contenido} & \textbf{Entregable sugerido} &
\textbf{Horas} \\
\midrule
1 & Introducción y primera mitad del Módulo~\ref{cap:fundamentos}. &
  Entorno instalado, datos generados y un script con tus primeros
  cálculos. & 7 \\
2 & Resto del Módulo~\ref{cap:fundamentos} e inicio del
  Módulo~\ref{cap:manipulacion}. &
  Un script que importa los CSV del curso y exporta una tabla a Excel. &
  8 \\
3 & Resto del Módulo~\ref{cap:manipulacion}. &
  Tabla de ventas por tienda y categoría, con clientes y productos
  unidos. & 8 \\
4 & Módulo~\ref{cap:visualizacion}. &
  Cinco gráficas de ventas con tu tema corporativo. & 8 \\
5 & Módulo~\ref{cap:eda}. &
  Informe exploratorio: distribuciones, atípicos y relaciones. & 8 \\
6 & Módulo~\ref{cap:bases-datos}. &
  Base SQLite con las tablas del curso y cinco consultas SQL. & 6 \\
7 & Módulo~\ref{cap:shiny}. &
  Dashboard Shiny de ventas con filtros. & 8 \\
8 & Módulo~\ref{cap:ml}. &
  Segmentación de clientes y pronóstico de ventas. & 10 \\
9 & Módulo~\ref{cap:reportes} y planeación del proyecto. &
  Reporte mensual automático en HTML y Word. & 7 \\
10 & Módulo~\ref{cap:proyecto}. &
  Proyecto final completo: ETL, análisis, dashboard y reporte. & 10 \\
\midrule
\multicolumn{3}{@{}l}{\textbf{Total aproximado}} & \textbf{80} \\
\bottomrule
\end{tabularx}
\end{table}

\begin{consejo}
Si tienes prisa por una necesidad concreta, estas rutas cortas funcionan:
para \textbf{reportes y gráficas}, Módulos~\ref{cap:fundamentos},
\ref{cap:manipulacion}, \ref{cap:visualizacion} y \ref{cap:reportes}; para
\textbf{dashboards}, añade el Módulo~\ref{cap:shiny}; para
\textbf{modelos predictivos}, añade los Módulos~\ref{cap:eda} y
\ref{cap:ml}. Los Módulos~\ref{cap:fundamentos} y \ref{cap:manipulacion} son
obligatorios en cualquier ruta: todo lo demás se construye sobre ellos.
\end{consejo}

% ----------------------------------------------------------------------------
\section{Convenciones del libro}
\label{sec:m0-convenciones}
% ----------------------------------------------------------------------------

Antes de empezar conviene que sepas leer el libro: cómo se distingue el
código que debes ejecutar, cómo se ve lo que R responde y para qué sirve
cada caja de color. Si todavía no tienes R instalado, solo lee esta sección;
en la sección~\ref{sec:m0-entorno} lo instalas y luego puedes volver a
probar los ejemplos.

\subsection{Bloques de código}

El libro usa cinco tipos de bloque. El más importante es el de
\textbf{código R que se ejecuta}: fondo gris claro y números de línea a la
izquierda. Todo el código de estos bloques se probó de principio a fin, en
orden, antes de publicar el libro.

\begin{rcode}
# Ventas diarias de una sucursal durante una semana (en pesos)
ventas_semana <- c(1200, 950, 1830, 1100, 2400, 3100, 2800)
sum(ventas_semana)    # venta total de la semana
mean(ventas_semana)   # venta promedio por día
\end{rcode}

Justo debajo aparece la \textbf{salida}: lo que R imprime en la consola al
ejecutar ese código. Tiene fondo blanco y una franja azul verdosa a la
izquierda:

\begin{rsalida}
[1] 13380
[1] 1911.429
\end{rsalida}

Los \textbf{números de línea} (1, 2, 3, 4 en gris) no forman parte del
código: sirven para que el texto pueda decir ``en la línea 2 creamos el
vector'' y tú sepas exactamente dónde mirar. No los escribas. Las líneas que
empiezan con \codigo{\#} son \termino{comentarios}: R las ignora; están ahí
para explicarte qué hace cada instrucción. En tus propios scripts, comenta
igual. Si una línea es demasiado larga para el ancho de la página, el libro
la parte y marca la continuación con el símbolo $\hookrightarrow$; en R
escríbela en una sola línea.

El segundo tipo es el \textbf{código R que no se ejecuta automáticamente}.
Se ve igual, pero con una franja naranja a la izquierda:

\begin{rnoejecutar}
# Instala un paquete (se hace UNA sola vez en cada computadora)
install.packages("tidyverse")
\end{rnoejecutar}

La franja naranja significa ``ejecútalo tú cuando lo necesites'': son
instalaciones, aplicaciones Shiny que abren una ventana, conexiones a
servidores que no existen en tu computadora, plantillas con valores
ficticios que debes cambiar o código con errores intencionales para
aprender a leerlos. Por eso casi nunca llevan salida debajo.

El tercer tipo es \textbf{SQL}, el lenguaje de las bases de datos, con una
franja morada. Lo aprenderás en el Módulo~\ref{cap:bases-datos}:

\begin{sqlcode}
-- Ventas totales por tienda, de mayor a menor
SELECT tienda_id, SUM(total) AS ventas_totales
FROM ventas
GROUP BY tienda_id
ORDER BY ventas_totales DESC;
\end{sqlcode}

Por último, los bloques de \textbf{texto literal} muestran cosas que no son
código R: encabezados de R Markdown, comandos de terminal, estructuras de
carpetas o archivos de configuración. Este, por ejemplo, es el encabezado
de un reporte que harás en el Módulo~\ref{cap:reportes}:

\begin{textocodigo}
---
title: "Reporte mensual de ventas"
author: "Área de Inteligencia Comercial"
output: html_document
---
\end{textocodigo}

\begin{advertencia}
No copies y pegues código desde el PDF: se copian también los números de
línea y a veces las comillas o los espacios cambian, lo que produce errores
difíciles de ver. Escribe el código tú mismo (así aprendes más rápido) o
abre el script del módulo en la carpeta \archivo{latex/codigo/}
(sección~\ref{sec:m0-codigo}).
\end{advertencia}

\subsection{Cómo leer una salida de R}
\label{sec:m0-leer-salida}

En la consola de RStudio verás el símbolo \codigo{>} antes de cada
instrucción: es el \termino{prompt}, la forma en que R dice ``estoy listo,
escribe algo''. En el libro no lo mostramos: el bloque de código contiene
exactamente lo que escribes y el bloque de salida, exactamente lo que R
responde.

\subsubsection{El \texorpdfstring{\codigo{[1]}}{[1]} y los índices}

Cada línea de salida de un vector empieza con un número entre corchetes. No
es un dato: es la \textbf{posición} del primer elemento de esa línea. Se nota
cuando el resultado no cabe en una sola línea:

\begin{rcode}
# Números de día de un mes de 30 días
dias_mes <- 1:30
dias_mes
\end{rcode}

\begin{rsalida}
 [1]  1  2  3  4  5  6  7  8  9 10 11 12 13 14 15 16 17 18 19 20 21 22
[23] 23 24 25 26 27 28 29 30
\end{rsalida}

La primera línea empieza en el elemento 1 y la segunda en el elemento 23.
En la salida de \codigo{sum(ventas\_semana)}, el \codigo{[1]} solo indica
que el resultado tiene un único elemento: 13380.

\subsubsection{Tablas: data frames y tibbles}

Los datos de negocio viven en tablas. En R, una tabla se llama
\termino{data frame}; el tidyverse usa una versión moderna llamada
\termino{tibble}, que se imprime de forma más informativa:

\begin{rcode}
library(tibble)   # paquete de tablas modernas (parte del tidyverse)

# Una tabla pequeña: ventas de tres tiendas
ventas_tienda <- tibble(
  tienda = c("Centro", "Norte", "Sur"),
  ventas = c(152300, 98750, 120400)
)
ventas_tienda
\end{rcode}

\begin{rsalida}
# A tibble: 3 x 2
  tienda ventas
  <chr>   <dbl>
1 Centro 152300
2 Norte   98750
3 Sur    120400
\end{rsalida}

Así se lee:
\begin{itemize}
  \item \codigo{\# A tibble: 3 x 2}: la tabla tiene 3 filas y 2 columnas.
        En tu pantalla verás el signo $\times$ en lugar de la \codigo{x}
        y algunos elementos en color; en el libro usamos caracteres
        simples.
  \item La segunda línea tiene los nombres de las columnas y la tercera su
        \textbf{tipo} entre \codigo{< >} (tabla~\ref{tab:m0-tipos}).
  \item Los números 1, 2, 3 de la izquierda son solo el número de fila.
  \item Si la tabla es grande, R muestra las primeras 10 filas y avisa
        con una línea como \codigo{\# i 355 more rows} (quedan 355 filas
        más) o \codigo{\# i 9 more variables: \ldots} (hay 9 columnas que no
        cupieron). En tu pantalla la \codigo{i} aparece como un símbolo de
        información. Los datos están completos: solo se recorta la
        impresión.
  \item Un número que termina en punto, como \codigo{126274.}, tiene
        decimales que no se muestran para ahorrar espacio.
\end{itemize}

\begin{table}[H]
\centering\small
\caption{Tipos de columna que verás en las salidas de R.}
\label{tab:m0-tipos}
\begin{tabularx}{\textwidth}{@{}l l X@{}}
\toprule
\textbf{Así aparece} & \textbf{Tipo} & \textbf{Ejemplo en los datos del
curso} \\
\midrule
\codigo{<chr>}  & Texto (\textit{character}) & nombre de la ciudad, método
  de pago \\
\codigo{<dbl>}  & Número decimal (\textit{double}) & total de la venta,
  salario \\
\codigo{<int>}  & Número entero (\textit{integer}) & conteo de ventas \\
\codigo{<lgl>}  & Lógico: \codigo{TRUE} o \codigo{FALSE} & ¿el cliente es
  premium? \\
\codigo{<date>} & Fecha & fecha de la venta \\
\codigo{<time>} & Hora del día & hora de la transacción \\
\codigo{<dttm>} & Fecha y hora & momento exacto de un registro \\
\codigo{<fct>}  & Factor: categoría con niveles fijos & segmento de
  cliente ordenado \\
\bottomrule
\end{tabularx}
\end{table}

\subsubsection{Valores faltantes, advertencias y errores}

En los datos reales faltan cosas. R marca un valor faltante con
\codigo{NA} (\textit{not available}). Mira qué pasa cuando una columna de
ventas llega como texto y uno de los valores dice ``N/D'':

\begin{rcode}
# Ventas que llegaron como texto desde un sistema
ventas_texto <- c("1500", "2300", "N/D", "980")
ventas_numero <- as.numeric(ventas_texto)   # convertir a número
\end{rcode}

R hace la conversión, pero responde en la consola con un aviso:

\begin{rsalida}
Warning message:
NAs introducidos por coerción
\end{rsalida}

Veamos qué quedó y qué pasa al sumar:

\begin{rcode}
ventas_numero
sum(ventas_numero)                 # la suma con un NA da NA
sum(ventas_numero, na.rm = TRUE)   # na.rm = TRUE ignora los NA
\end{rcode}

\begin{rsalida}
[1] 1500 2300   NA  980
[1] NA
[1] 4780
\end{rsalida}

El texto ``N/D'' se convirtió en \codigo{NA}; la suma ``normal'' devuelve
\codigo{NA} (si falta un dato, R no puede saber el total), y con
\codigo{na.rm = TRUE} obtienes la suma de los valores disponibles, 4780.

Hay tres clases de avisos. Una \termino{advertencia}
(\textit{warning}) significa ``hice lo que pediste, pero algo raro pasó'':
aquí R no pudo convertir ``N/D'' en número y puso \codigo{NA}. Un
\termino{error} (\textit{error}) significa ``me detuve y no hice nada''. Un
\termino{mensaje} (\textit{message}) es solo informativo. Nunca ignores una
advertencia sin entenderla: en BI suele ser la primera pista de un problema
de calidad de datos.

Fíjate también en que el aviso salió mitad en inglés (\textit{Warning
message}) y mitad en español. R traduce muchos de sus mensajes según el
idioma de tu sistema, así que en tu computadora pueden salir en español, en
inglés o mezclados. No te preocupes: el significado es el mismo. Por último,
cuando una salida es muy larga la recortamos y lo indicamos con una línea
\codigo{...}.

\subsection{Las cajas didácticas}

Además del código, el libro usa cajas de colores. Cada una tiene un
propósito fijo; al verla sabrás qué tipo de contenido esperar. Al inicio de
cada módulo, una caja azul de \textbf{objetivos} (como la de la primera
página de esta introducción) lista lo que podrás hacer al terminarlo. Estas
son las demás.

La caja \textbf{Concepto clave} contiene definiciones e ideas teóricas que
necesitas para seguir adelante. Si vas a repasar un módulo en diez minutos,
lee solo estas cajas:

\begin{concepto}{ticket promedio}
El \termino{ticket promedio} es el importe medio de cada compra:
ventas totales divididas entre el número de compras (tickets). Si sube, tus
clientes gastan más por visita; si baja, compran menos cosas o más baratas.
\end{concepto}

La caja \textbf{Ejemplo práctico} plantea un problema de negocio que se
resuelve paso a paso. Normalmente contiene la pregunta y el código va justo
después:

\begin{ejemplo}{ticket promedio de la semana}
Con las ventas de la semana (\codigo{ventas\_semana}) y el número de
tickets de cada día, ¿cuál fue el ticket promedio de la sucursal?
\end{ejemplo}

\begin{rcode}
# Número de tickets (compras) registrados cada día de la semana
tickets_semana <- c(40, 31, 55, 38, 70, 88, 81)

# Ticket promedio = venta total / número total de tickets
ticket_promedio <- sum(ventas_semana) / sum(tickets_semana)
round(ticket_promedio, 2)   # redondeamos a centavos
\end{rcode}

\begin{rsalida}
[1] 33.2
\end{rsalida}

Interpretación: en esa semana cada cliente gastó en promedio \$33.20 por
compra. Después de cada salida importante encontrarás una interpretación
así: el número por sí solo no sirve de nada hasta que dices qué significa
para el negocio.

La caja \textbf{En la empresa} cuenta cómo se usa una técnica en una
organización real, con sus ventajas y sus problemas prácticos:

\begin{negocio}{el reporte de los lunes}
En muchas cadenas, cada lunes alguien arma a mano el reporte de ventas de
la semana anterior: descarga archivos, los pega en Excel y actualiza
gráficas. Con un script de R y un reporte de R Markdown
(Módulo~\ref{cap:reportes}) ese trabajo se reduce a ejecutar un archivo, y
el analista usa ese tiempo en lo que sí aporta valor: explicar qué pasó y
qué hacer.
\end{negocio}

La caja \textbf{Consejo} reúne buenas prácticas profesionales, esas que no
son obligatorias para que el código funcione pero que distinguen a un
analista cuidadoso:

\begin{consejo}
Nombra tus scripts con un número y un verbo que indiquen el orden en que se
ejecutan: \archivo{01\_importar.R}, \archivo{02\_limpiar.R},
\archivo{03\_analizar.R}. Dentro de un año agradecerás saber por dónde
empezar.
\end{consejo}

La caja \textbf{Cuidado} advierte sobre errores comunes y cómo evitarlos:

\begin{advertencia}
R distingue entre mayúsculas y minúsculas: \codigo{Ventas} y
\codigo{ventas} son dos objetos distintos. Si creas \codigo{ventas} y luego
escribes \codigo{Ventas}, R responderá con un error que dice que el objeto
no se encontró (\textit{object not found}).
\end{advertencia}

La caja \textbf{Hazlo tú} es un reto corto para practicar en ese mismo
momento. Muchas traen una pista, pero no una solución formal: la idea es
que experimentes.

\begin{hazlotu}
Cambia los valores de \codigo{ventas\_semana} por las ventas de un negocio
que conozcas (o inventa unas) y vuelve a calcular el total, el promedio
diario y el ticket promedio.
\end{hazlotu}

Al final de cada módulo, la caja \textbf{Resumen del módulo} repasa lo
aprendido y suele incluir una tabla con las funciones clave:

\begin{resumen}[title={Resumen del módulo (ejemplo)}]
\begin{itemize}
  \item \codigo{sum()} suma y \codigo{mean()} promedia los elementos de un
        vector.
  \item \codigo{NA} marca un valor faltante; \codigo{na.rm = TRUE} lo
        ignora en los cálculos.
\end{itemize}
\end{resumen}

Después del resumen vienen los \textbf{ejercicios}, en cajas naranjas
numeradas (Ejercicio 2.1, 2.2\ldots). Cada ejercicio indica su nivel
(Básico, Intermedio, Avanzado o Reto) y, en la esquina inferior derecha, la
página del Apéndice~\ref{ap:soluciones} donde está su solución completa:
explicación, código, salida real e interpretación. Intenta resolverlo
antes de mirar la solución.

\subsection{Formato del texto}

Dentro de los párrafos, el formato también te da pistas:
\codigo{summarise()} es código de R (funciones, objetos, argumentos);
\paquete{dplyr}, el nombre de un paquete; \archivo{ventas\_retail.csv}, un
archivo o carpeta; \menu{Tools > Global Options}, un menú de RStudio;
\tecla{Ctrl} + \tecla{Enter}, un atajo de teclado (en Mac cambia
\tecla{Ctrl} por \tecla{Cmd}); y un \termino{término en azul} es un
concepto nuevo que se define ahí mismo.

% ----------------------------------------------------------------------------
\section{Preparar tu entorno de trabajo}
\label{sec:m0-entorno}
% ----------------------------------------------------------------------------

Vas a instalar dos programas, en este orden: primero \textbf{R} (el motor)
y después \textbf{RStudio} (el tablero). Luego configurarás RStudio,
instalarás los paquetes del curso y crearás un Proyecto de RStudio. Calcula
entre 30 y 60 minutos, la mayor parte esperando descargas. El libro se
verificó con R 4.3.3; cualquier versión 4.3 o posterior sirve.

\subsection{Paso 1: instalar R}

R se descarga de \termino{CRAN} (\textit{Comprehensive R Archive
Network}), el repositorio oficial: \url{https://cran.r-project.org}.

\textbf{Windows}
\begin{enumerate}
  \item En la página de CRAN elige \menu{Download R for Windows} y luego
        \menu{base}.
  \item Pulsa \menu{Download R-4.x.x for Windows} (el número cambia con
        cada versión).
  \item Ejecuta el instalador \archivo{.exe} y acepta las opciones por
        defecto.
  \item Verás una página sobre \textbf{Rtools}: solo hace falta para
        compilar paquetes desde su código fuente; para este curso no la
        necesitas.
\end{enumerate}

\textbf{macOS}
\begin{enumerate}
  \item En CRAN elige \menu{Download R for macOS}.
  \item Descarga el instalador \archivo{.pkg} que corresponde a tu
        procesador: \textbf{Apple silicon} (M1, M2, M3\ldots; archivo con
        \archivo{arm64} en el nombre) o \textbf{Intel} (\archivo{x86\_64}).
        Si no sabes cuál tienes, abre el menú Apple $\rightarrow$
        \menu{Acerca de esta Mac} y mira el campo ``Chip'' o
        ``Procesador''.
  \item Abre el \archivo{.pkg} y sigue el asistente.
\end{enumerate}

\textbf{Linux (Ubuntu o Debian)}. Abre una terminal y ejecuta:

\begin{textocodigo}
sudo apt update
sudo apt install r-base r-base-dev
# Bibliotecas del sistema que necesitan algunos paquetes de R
sudo apt install libcurl4-openssl-dev libssl-dev libxml2-dev
R --version
\end{textocodigo}

La versión de los repositorios de tu distribución puede estar algo
atrasada; en la sección de Linux de CRAN hay instrucciones para agregar su
repositorio y obtener siempre la versión más reciente.

\subsection{Paso 2: instalar RStudio}

\begin{enumerate}
  \item Entra a \url{https://posit.co/download/rstudio-desktop/}.
  \item Descarga \textbf{RStudio Desktop} (gratuito) para tu sistema. La
        página suele detectarlo sola.
  \item Instálalo. En Windows, ejecuta el archivo \archivo{.exe}; en
        macOS, abre el \archivo{.dmg} y arrastra RStudio a la carpeta
        \archivo{Aplicaciones}; en Ubuntu, instala el archivo
        \archivo{.deb}.
  \item Abre RStudio. Si en la consola aparece un texto que empieza con
        \codigo{R version 4.}, R y RStudio están conectados.
\end{enumerate}

\begin{consejo}
Si no puedes instalar software en la computadora de tu trabajo, crea una
cuenta gratuita en \textbf{Posit Cloud} (\url{https://posit.cloud}): es
RStudio en el navegador. Sube ahí la carpeta del curso y sigue el libro
igual. El plan gratuito tiene límites de memoria y de horas de uso, pero
alcanza para practicar.
\end{consejo}

\subsection{Un vistazo a RStudio}

Al abrir RStudio verás cuatro paneles: el \textbf{editor} (arriba a la
izquierda), donde escribes y guardas tus scripts; la \textbf{consola}
(abajo a la izquierda), donde R ejecuta el código y responde; el
\textbf{entorno} (arriba a la derecha), con los objetos que has creado; y
un panel con pestañas para archivos, gráficas, paquetes y ayuda (abajo a la
derecha). Si solo ves tres, abre un script nuevo con
\menu{File > New File > R Script}. El atajo que más usarás es
\tecla{Ctrl} + \tecla{Enter} (\tecla{Cmd} + \tecla{Return} en Mac): envía a
la consola la línea donde está el cursor, o el texto seleccionado. El
Módulo~\ref{cap:fundamentos} recorre RStudio con detalle y el
Apéndice~\ref{ap:referencia} reúne los atajos de teclado.

\subsection{Paso 3: configurar RStudio}

Abre \menu{Tools > Global Options} y ajusta las opciones de la
tabla~\ref{tab:m0-config}. La más importante es la primera pareja: evita
que R guarde y restaure automáticamente los objetos de la sesión anterior.

\begin{table}[H]
\centering\small
\caption{Configuración recomendada de RStudio
(\menu{Tools > Global Options}).}
\label{tab:m0-config}
\begin{tabularx}{\textwidth}{@{}>{\raggedright\arraybackslash}p{2.2cm}
    >{\raggedright\arraybackslash}p{4.6cm}
    >{\raggedright\arraybackslash}p{2.1cm} X@{}}
\toprule
\textbf{Sección} & \textbf{Opción} & \textbf{Valor} & \textbf{Por qué} \\
\midrule
\menu{General} & Restore .RData into workspace at startup & Desmarcada &
  Cada sesión empieza limpia, sin objetos ``fantasma'' de ayer. \\
\menu{General} & Save workspace to .RData on exit & Never &
  Tu trabajo debe vivir en scripts, no en la memoria. \\
\menu{Code > Editing} & Insert spaces for Tab (Tab width 2) & Marcada &
  Sangría uniforme, al estilo tidyverse. \\
\menu{Code > Editing} & Use native pipe operator, \codigo{|>} &
  Desmarcada & El libro usa \codigo{\%>\%}; así el atajo inserta ese pipe
  (el Módulo~\ref{cap:manipulacion} explica los dos). \\
\menu{Code > Display} & Show line numbers; Highlight selected line &
  Marcadas & Ubicas los errores por número de línea. \\
\menu{Code > Display} & Show margin; Use rainbow parentheses & Marcadas &
  Líneas cortas y paréntesis de colores para ver cuál cierra a cuál. \\
\menu{Code > Saving} & Default text encoding & UTF-8 &
  Acentos y eñes correctos en cualquier sistema. \\
\menu{Appearance} & Editor theme, font size & A tu gusto (12 a 14) &
  Trabajarás muchas horas: que sea cómodo. \\
\bottomrule
\end{tabularx}
\end{table}

\subsection{Paso 4: instalar los paquetes del curso}

Un \termino{paquete} se \textbf{instala} una sola vez en cada computadora
(con \codigo{install.packages()}) y se \textbf{carga} en cada sesión en que
lo uses (con \codigo{library()}). Es como comprar un libro y después
abrirlo cada vez que lo lees. Copia el siguiente bloque en un script nuevo y
ejecútalo completo; tarda entre 10 y 30 minutos según tu conexión.

\begin{rnoejecutar}[title=Paquetes obligatorios del curso, colbacktitle=primario!85, coltitle=white, fonttitle=\sffamily\bfseries\small]
paquetes <- c(
  # --- Importar, manipular, visualizar y explorar (Módulos 1 a 4) ---
  "tidyverse",      # dplyr, ggplot2, tidyr, readr, tibble, stringr...
  "readxl",         # leer archivos de Excel (.xlsx y .xls)
  "writexl",        # escribir archivos de Excel
  "lubridate",      # trabajar con fechas: meses, semanas, días
  "scales",         # formato de ejes: $, %, separador de miles
  "skimr",          # resumen rápido de una tabla completa
  "corrplot",       # gráficas de matrices de correlación
  # --- Bases de datos y SQL (Módulo 5) ---
  "DBI",            # interfaz común para cualquier base de datos
  "RSQLite",        # base de datos SQLite en un archivo local
  "dbplyr",         # traducir verbos de dplyr a SQL
  # --- Dashboards interactivos (Módulo 6) ---
  "shiny",          # aplicaciones web con R
  "shinydashboard", # plantilla de dashboard para Shiny
  "DT",             # tablas interactivas: buscar, ordenar, filtrar
  "plotly",         # gráficas interactivas
  # --- Machine learning (Módulo 7) ---
  "cluster",        # segmentación (clustering) y silueta
  "factoextra",     # visualizar clusters y componentes principales
  "forecast",       # series de tiempo y pronósticos
  "rpart",          # árboles de decisión
  "randomForest",   # bosques aleatorios
  "caret",          # entrenar y evaluar modelos
  # --- Reportes automáticos (Módulo 8) ---
  "rmarkdown",      # reportes en HTML, Word y PDF
  "knitr"           # motor que ejecuta el código de los reportes
)

# Instala solo los paquetes que todavía no tienes
faltantes <- setdiff(paquetes, rownames(installed.packages()))
install.packages(faltantes)
\end{rnoejecutar}

Algunos módulos mencionan paquetes adicionales. No son necesarios para
seguir el libro, pero puedes instalarlos si quieres probar esos ejemplos:

\begin{rnoejecutar}[title=Paquetes opcionales, colbacktitle=primario!85, coltitle=white, fonttitle=\sffamily\bfseries\small]
install.packages(c(
  "arules",      # reglas de asociación (análisis de canasta de compras)
  "rpart.plot",  # dibujos más legibles de árboles de decisión
  "janitor",     # limpiar nombres de columnas y tablas de frecuencia
  "openxlsx"     # Excel con formato: colores, anchos, estilos
))
\end{rnoejecutar}

\begin{advertencia}
Si durante la instalación R pregunta \textit{Do you want to install from
sources the packages which need compilation?}, responde \codigo{no}: se
instalará la versión ya compilada, que funciona igual y no requiere
herramientas adicionales. En Linux los paquetes siempre se compilan; si
alguno falla, casi siempre es porque falta una biblioteca del sistema (el
mensaje de error dice cuál) como las del paso~1.
\end{advertencia}

\subsection{Paso 5: crear un Proyecto de RStudio}

Descarga la carpeta del curso, \archivo{curso-bi-rstudio/} (por ejemplo,
descargando el repositorio como ZIP y descomprimiéndolo), y colócala en un
lugar fácil de encontrar, como \archivo{Documentos}. Evita rutas con
caracteres especiales o carpetas que se sincronizan constantemente con la
nube, porque a veces bloquean archivos mientras R los escribe. Después:

\begin{enumerate}
  \item En RStudio elige \menu{File > New Project\ldots}.
  \item Selecciona \menu{Existing Directory} y, con \menu{Browse\ldots},
        la carpeta \archivo{curso-bi-rstudio}.
  \item Pulsa \menu{Create Project}. RStudio se reinicia, crea el archivo
        \archivo{curso-bi-rstudio.Rproj} y muestra el nombre del proyecto en
        la esquina superior derecha.
  \item Las próximas veces abre el curso con doble clic en
        \archivo{curso-bi-rstudio.Rproj} (o desde \menu{File > Open
        Project}).
\end{enumerate}

\begin{concepto}{Proyecto de RStudio y rutas relativas}
Un \termino{Proyecto de RStudio} es una carpeta con un archivo
\archivo{.Rproj}. Al abrirlo, RStudio fija esa carpeta como
\termino{directorio de trabajo}: el lugar desde donde R busca y guarda
archivos. Gracias a eso puedes usar \termino{rutas relativas} como
\codigo{"datasets/ventas\_retail.csv"} en lugar de rutas absolutas como
\codigo{"C:/Users/Ana/Documents/\ldots"}. Tu código funcionará igual en
cualquier computadora donde copies la carpeta, y cada proyecto tendrá su
propio historial y sus propios archivos abiertos.
\end{concepto}

Para comprobar dónde está trabajando R, escribe en la consola:

\begin{rnoejecutar}
getwd()   # muestra el directorio de trabajo actual
\end{rnoejecutar}

La respuesta debe terminar en \archivo{curso-bi-rstudio}. Todo el código
del libro supone que el directorio de trabajo es esa carpeta.

\begin{advertencia}
Evita escribir \codigo{setwd("C:/Users/TuNombre/\ldots")} al inicio de tus
scripts. Funciona en tu computadora, pero falla en la de cualquier otra
persona (y en la tuya, si mueves la carpeta). Con un Proyecto de RStudio no
lo necesitas.
\end{advertencia}

\subsection{La estructura de carpetas del curso}

La carpeta del curso está organizada así:

\begin{textocodigo}
curso-bi-rstudio/
|-- curso-bi-rstudio.Rproj    <- lo crea RStudio (paso 5)
|-- README.md                 <- presentación del curso
|-- GUIA_INSTALACION.md       <- guía de instalación resumida
|-- datasets/                 <- datos del curso
|   |-- generar_datasets.R    <- script que crea todos los datos
|   |-- ventas_retail.csv     <- (estos archivos aparecen al
|   |-- clientes.csv          <-  ejecutar el script)
|   |-- ...
|   `-- datos_empresa.xlsx
|-- latex/                    <- el libro que estás leyendo
|   |-- curso-bi-rstudio.tex  <- archivo principal del libro
|   |-- preambulo.tex         <- estilos, colores y cajas
|   |-- capitulos/            <- un archivo .tex por módulo
|   |-- apendices/            <- soluciones y referencia rápida
|   |-- figuras/              <- gráficas del libro en PDF
|   |-- codigo/               <- un script .R por módulo
|   `-- herramientas/         <- extraer código, generar figuras
|-- modulo-01-fundamentos/    <- scripts complementarios de
|-- modulo-02-.../            <-  práctica de algunos módulos
|-- proyecto-final/           <- material del proyecto integrador
`-- resultados/               <- lo crean los ejemplos que guardan
                                 archivos (CSV, Excel, gráficas)
\end{textocodigo}

Las carpetas \archivo{modulo-XX/} contienen scripts de práctica
complementarios; el código completo y verificado de cada módulo del libro
está en \archivo{latex/codigo/}.

% ----------------------------------------------------------------------------
\section{Generar los datos del curso}
\label{sec:m0-generar}
% ----------------------------------------------------------------------------

Los datos del curso no se descargan: los crea un script de R. Así pesan
poco, y todos obtenemos exactamente los mismos datos.

\subsection{Paso 6: ejecutar el generador}

Con el proyecto abierto (el directorio de trabajo debe ser
\archivo{curso-bi-rstudio}), escribe en la consola:

\begin{rcode}
source("datasets/generar_datasets.R")   # ejecuta el script completo
\end{rcode}

\begin{rsalida}
[1] "✓ ventas_retail.csv creado"
[1] "✓ clientes.csv creado"
[1] "✓ productos.csv creado"
[1] "✓ empleados.csv creado"
[1] "✓ transacciones.csv creado"
[1] "✓ tiendas.csv creado"
[1] "✓ datos_empresa.xlsx creado"

============================================
  DATASETS GENERADOS EXITOSAMENTE
============================================
Archivos creados:
  1. ventas_retail.csv ( 365 filas )
  2. clientes.csv ( 200 filas )
  3. productos.csv ( 50 filas )
  4. empleados.csv ( 100 filas )
  5. transacciones.csv ( 1000 filas )
  6. tiendas.csv ( 10 filas )
  7. datos_empresa.xlsx (multi-hoja)
============================================
¡Listo para usar en los módulos del curso!
============================================
\end{rsalida}

La función \codigo{source()} ejecuta de principio a fin todas las
instrucciones de un script. El generador fija una \termino{semilla}
aleatoria con \codigo{set.seed(123)}: los datos son aleatorios, pero la
semilla hace que R produzca siempre la misma secuencia de números. Por eso,
con R 3.6 o posterior, tus archivos serán idénticos a los del libro y tus
resultados coincidirán con las salidas impresas. Solo necesitas ejecutarlo
una vez; si borras algún archivo por accidente, vuelve a ejecutarlo.

\begin{advertencia}
Si aparece \textit{cannot open file 'datasets/generar\_datasets.R': No such
file or directory}, R no está trabajando en la carpeta del curso. Abre el
archivo \archivo{curso-bi-rstudio.Rproj} y vuelve a intentarlo. Si el error
menciona \paquete{writexl}, te falta instalar los paquetes del paso~4.
\end{advertencia}

\subsection{Nota sobre el idioma de los días y los meses}

El generador escribe el día de la semana con la función
\codigo{weekdays()}, que usa el idioma configurado en tu sistema para las
fechas. Puedes consultarlo así:

\begin{rcode}
Sys.getlocale("LC_TIME")   # idioma y región que R usa para las fechas
\end{rcode}

\begin{rsalida}
[1] "es_MX.UTF-8"
\end{rsalida}

Los datos del libro se generaron con español de México
(\codigo{es\_MX.UTF-8}), por eso verás ``lunes'', ``martes''\ldots{} En
Windows la respuesta se parece a \codigo{"Spanish\_Mexico.utf8"}. Si en tu
computadora ves algo como \codigo{"English\_United States.utf8"} o
\codigo{"en\_US.UTF-8"}, tus días saldrán en inglés (``Monday'',
``Tuesday''\ldots) y algunas salidas del libro no coincidirán con las tuyas.
Para generar los datos en español, cambia el idioma de las fechas antes de
ejecutar el generador:

\begin{rnoejecutar}
# Windows
Sys.setlocale("LC_TIME", "Spanish")
# macOS y Linux
Sys.setlocale("LC_TIME", "es_MX.UTF-8")

source("datasets/generar_datasets.R")   # vuelve a generar los datos
\end{rnoejecutar}

El cambio dura solo esa sesión de R. En Linux, si la instrucción responde
con una advertencia, el idioma no está instalado; se agrega con
\codigo{sudo locale-gen es\_MX.UTF-8} en la terminal. Las cifras, los
identificadores y los totales no dependen del idioma: solo los nombres de
días y meses.

\subsection{Paso 7: verificar la instalación}

El siguiente script revisa todo de una vez: la versión de R, que cada
paquete del curso se pueda cargar y que existan los archivos de datos. Es
un buen primer script para guardar en tu proyecto como
\archivo{00\_verificar\_entorno.R}. Lee los comentarios: aunque todavía no
conozcas \codigo{for} ni \codigo{if}, entenderás la idea.

Primero carga el tidyverse, que el script usa para leer los CSV:

\begin{rcode}
library(tidyverse)   # carga dplyr, ggplot2, readr y compañía
\end{rcode}

Al ejecutar \codigo{library(tidyverse)} verás un mensaje de bienvenida con
la lista de paquetes que se cargaron y un par de avisos de
``conflictos''. No es un error y no lo reproducimos en el libro (el
Módulo~\ref{cap:fundamentos} lo explica). Ahora sí, el script de
verificación:

\begin{rcode}[title=Verificación del entorno del curso, colbacktitle=primario!85, coltitle=white, fonttitle=\sffamily\bfseries\small]
# ---- 1. Versión de R ----
cat("Versión de R:", as.character(getRversion()), "\n\n")

# ---- 2. Paquetes del curso: ¿se pueden cargar? ----
paquetes_curso <- c("tidyverse", "readxl", "writexl", "lubridate",
                    "scales", "DBI", "RSQLite", "dbplyr", "shiny",
                    "shinydashboard", "DT", "plotly", "cluster",
                    "factoextra", "forecast", "rpart", "randomForest",
                    "caret", "rmarkdown", "knitr", "skimr", "corrplot")
problemas <- 0                         # contador de problemas

cat("Paquetes:\n")
for (p in paquetes_curso) {            # revisa los paquetes uno por uno
  # ¿carga bien? (suppressMessages oculta avisos técnicos de carga)
  if (suppressMessages(requireNamespace(p, quietly = TRUE))) {
    version <- as.character(packageVersion(p))
    cat(sprintf("  OK     %-15s %s\n", p, version))
  } else {
    problemas <- problemas + 1
    cat(sprintf("  FALTA  %-15s instálalo (paso 4)\n", p))
  }
}

# ---- 3. Archivos de datos ----
cat("\nArchivos en datasets/:\n")
tablas <- c("ventas_retail", "clientes", "productos", "empleados",
            "transacciones", "tiendas")
for (t in tablas) {
  archivo_csv <- file.path("datasets", paste0(t, ".csv"))
  if (file.exists(archivo_csv)) {
    datos <- read_csv(archivo_csv, show_col_types = FALSE)
    cat(sprintf("  OK     %-18s %5d filas %3d columnas\n",
                basename(archivo_csv), nrow(datos), ncol(datos)))
  } else {
    problemas <- problemas + 1
    cat(sprintf("  FALTA  %s (ejecuta el paso 6)\n", archivo_csv))
  }
}
hojas <- readxl::excel_sheets("datasets/datos_empresa.xlsx")
cat("  OK     datos_empresa.xlsx con hojas:\n        ",
    paste(hojas, collapse = ", "), "\n")

# ---- 4. Veredicto ----
if (problemas == 0) {
  cat("\n¡Todo listo! Tu entorno está preparado para el curso.\n")
} else {
  cat("\nHay", problemas, "problema(s): revisa las líneas FALTA.\n")
}
\end{rcode}

\begin{rsalida}
Versión de R: 4.3.3

Paquetes:
  OK     tidyverse       2.0.0
  OK     readxl          1.4.3
  OK     writexl         1.5.0
  OK     lubridate       1.9.3
...
  OK     skimr           2.1.5
  OK     corrplot        0.92

Archivos en datasets/:
  OK     ventas_retail.csv    365 filas  14 columnas
  OK     clientes.csv         200 filas  10 columnas
  OK     productos.csv         50 filas  13 columnas
  OK     empleados.csv        100 filas  12 columnas
  OK     transacciones.csv   1000 filas  14 columnas
  OK     tiendas.csv           10 filas   8 columnas
  OK     datos_empresa.xlsx con hojas:
         Ventas, Clientes, Productos, Empleados, Tiendas

¡Todo listo! Tu entorno está preparado para el curso.
\end{rsalida}

Recortamos la lista de paquetes (la línea \codigo{...}); en tu pantalla
aparecen los 22. Tus números de versión serán
distintos si instalaste paquetes más recientes: no importa. Lo que importa
es que no haya ninguna línea \codigo{FALTA}.

\subsection{Tu primer análisis}

Ya que todo funciona, date un gusto: responde tu primera pregunta de
negocio. ¿Cuánto vendió la cadena en cada trimestre de 2023? No te
preocupes por entender cada función; en el Módulo~\ref{cap:manipulacion}
las verás con calma. Lee el código como una receta, gracias al operador
\codigo{\%>\%} (el \termino{pipe}, que se lee ``y luego''):

\begin{rcode}
# Leer la tabla de ventas
ventas <- read_csv("datasets/ventas_retail.csv", show_col_types = FALSE)

ventas %>%                                  # toma las ventas, y luego
  group_by(trimestre) %>%                   # agrúpalas por trimestre,
  summarise(ventas_totales = sum(total),    # suma el total vendido,
            numero_ventas = n(),            # cuenta las ventas
            ticket_promedio = mean(total),  # y promedia su importe
            .groups = "drop")
\end{rcode}

\begin{rsalida}
# A tibble: 4 x 4
  trimestre ventas_totales numero_ventas ticket_promedio
  <chr>              <dbl>         <int>           <dbl>
1 Q1               126274.            90           1403.
2 Q2               120717.            91           1327.
3 Q3               134427.            92           1461.
4 Q4               132498.            92           1440.
\end{rsalida}

¿Qué nos dice esto? El tercer trimestre fue el mejor, con unos \$134,427, y
el segundo el más flojo, con \$120,717. Como el número de ventas casi no
cambia entre trimestres (de 90 a 92), la diferencia viene del
\textbf{ticket promedio}: en el tercer trimestre cada venta fue de \$1,461
en promedio, frente a \$1,327 en el segundo. Ese es exactamente el tipo de
lectura que harás en todo el libro: pasar de la tabla a la frase que le
sirve a quien decide.

\begin{hazlotu}
Cambia \codigo{group\_by(trimestre)} por \codigo{group\_by(mes)} y vuelve a
ejecutar. ¿Qué mes vendió más? Luego prueba con
\codigo{group\_by(dia\_semana)}.
\end{hazlotu}

% ----------------------------------------------------------------------------
\section{Diccionario de datos}
\label{sec:m0-diccionario}
% ----------------------------------------------------------------------------

Un \termino{diccionario de datos} describe cada tabla y cada columna: qué
significa, de qué tipo es y qué valores puede tomar. Es el primer
documento que pide un analista al llegar a una empresa, y el que más tiempo
ahorra. Vuelve a esta sección siempre que tengas dudas sobre una columna.

\subsection{La empresa y sus tablas}

Los datos describen el año 2023 de una cadena minorista ficticia con
\textbf{10 tiendas} en 6 ciudades (CDMX, Guadalajara, Monterrey, Puebla,
Querétaro y Tijuana), \textbf{50 productos} de 10 categorías,
\textbf{200 clientes} y \textbf{100 empleados}. Las tablas siguen la
organización típica de un almacén de datos de BI.

\begin{concepto}{tablas de hechos y dimensiones (esquema estrella)}
Una \termino{tabla de hechos} registra eventos del negocio con medidas
numéricas: cada venta, cada transacción. Una \termino{tabla de dimensión}
describe el contexto de esos eventos: quién compró (clientes), qué se
vendió (productos), dónde (tiendas) y quién atendió (empleados). Las tablas
se conectan mediante \termino{llaves}: la \termino{llave primaria} (PK)
identifica cada fila de una dimensión, y la \termino{llave foránea} (FK) de
la tabla de hechos apunta a ella. Dibujadas, las dimensiones rodean a los
hechos como las puntas de una estrella: por eso se llama
\termino{esquema estrella}.
\end{concepto}

\begin{figure}[H]
\centering
\begin{tikzpicture}[
  cab/.style={draw=#1, fill=#1, text=white, font=\sffamily\bfseries\small,
    minimum width=3.7cm, minimum height=0.55cm, inner sep=2pt},
  cuerpo/.style={draw=#1, fill=#1!5, font=\ttfamily\scriptsize,
    text width=3.4cm, inner sep=0.15cm, align=left},
  rel/.style={line width=0.9pt, draw=#1},
  llave/.style={font=\ttfamily\scriptsize, text=#1, fill=white,
    inner sep=1pt},
]
  % ---- Tablas de hechos (centro) ----
  \node[cab=acento] (vh) at (0,3.3) {ventas\_retail};
  \node[cuerpo=acento, below=0pt of vh] (vb)
    {fecha\\
     \textbf{producto\_id} (FK)\\
     \textbf{cliente\_id} (FK)\\
     \textbf{tienda\_id} (FK)\\
     \textbf{vendedor\_id} (FK)\\
     cantidad, total, \ldots};
  \node[cab=acento] (th) at (0,-0.35) {transacciones};
  \node[cuerpo=acento, below=0pt of th] (tb)
    {\textbf{transaccion\_id} (PK)\\
     \textbf{producto\_id} (FK)\\
     \textbf{cliente\_id} (FK)\\
     \textbf{tienda\_id} (FK)\\
     \textbf{vendedor\_id} (FK)\\
     fecha, hora, total\_\ldots};
  % ---- Dimensiones (izquierda) ----
  \node[cab=secundario] (ch) at (-6.2,3.3) {clientes};
  \node[cuerpo=secundario, below=0pt of ch] (cb)
    {\textbf{cliente\_id} (PK)\\
     nombre, edad, ciudad\\
     segmento, \ldots};
  \node[cab=morado] (ph) at (-6.2,-0.35) {productos};
  \node[cuerpo=morado, below=0pt of ph] (pb)
    {\textbf{producto\_id} (PK)\\
     categoria, marca\\
     precio\_catalogo, costo, \ldots};
  % ---- Dimensiones (derecha) ----
  \node[cab=exito] (sh) at (6.2,3.3) {tiendas};
  \node[cuerpo=exito, below=0pt of sh] (sb)
    {\textbf{tienda\_id} (PK)\\
     nombre\_tienda, ciudad\\
     tipo, tamano\_m2, \ldots};
  \node[cab=primario] (eh) at (6.2,-0.35) {empleados};
  \node[cuerpo=primario, below=0pt of eh] (eb)
    {\textbf{empleado\_id} (PK)\\
     nombre, departamento\\
     puesto, salario, \ldots};
  % ---- Relaciones ----
  \draw[rel=secundario] (vb.west) -- (cb.east);
  \draw[rel=secundario] (tb.west) -- (cb.east);
  \draw[rel=morado] (vb.west) -- (pb.east);
  \draw[rel=morado] (tb.west) -- (pb.east);
  \draw[rel=exito] (vb.east) -- (sb.west);
  \draw[rel=exito] (tb.east) -- (sb.west);
  \draw[rel=primario] (vb.east) -- (eb.west);
  \draw[rel=primario] (tb.east) -- (eb.west);
  % Etiquetas de las llaves junto a cada dimensión
  \node[llave=secundario, anchor=south west] at ([xshift=2pt,yshift=3pt]cb.east)
    {cliente\_id};
  \node[llave=morado, anchor=north west] at ([xshift=2pt,yshift=-3pt]pb.east)
    {producto\_id};
  \node[llave=exito, anchor=south east] at ([xshift=-2pt,yshift=3pt]sb.west)
    {tienda\_id};
  \node[llave=primario, anchor=north east, align=right]
    at ([xshift=-2pt,yshift=-3pt]eb.west)
    {vendedor\_id\\= empleado\_id};
\end{tikzpicture}
\caption{Relaciones entre las tablas del curso. Las dos tablas de hechos
(al centro) se unen con las cuatro dimensiones mediante llaves foráneas
(FK) que apuntan a la llave primaria (PK) de cada dimensión. Cada fila de
una dimensión puede aparecer en muchas ventas (relación de muchos a uno).}
\label{fig:m0-relaciones}
\end{figure}

Observa la relación más delicada: en las tablas de hechos la columna se
llama \codigo{vendedor\_id}, pero en \archivo{empleados.csv} se llama
\codigo{empleado\_id}. Cuando unas estas tablas en el
Módulo~\ref{cap:manipulacion} tendrás que indicarle a R que son la misma
llave con nombres distintos.

En las tablas siguientes, el \textbf{tipo} es el que asigna
\codigo{read\_csv()} al leer cada archivo (tabla~\ref{tab:m0-tipos}). Los
identificadores y las cantidades se leen como \codigo{<dbl>} aunque sean
enteros; es normal.

\subsection{\texorpdfstring{\archivo{ventas\_retail.csv}}{ventas\_retail.csv}:
ventas diarias}

\textbf{365 filas y 14 columnas}. Cada fila es una venta y hay
exactamente una venta por cada día de 2023. Es una simplificación: una
cadena real tiene cientos de tickets diarios. Sirve para practicar
análisis por fecha, tienda, producto y descuento.

\begin{table}[H]
\centering\small
\caption{Diccionario de \archivo{ventas\_retail.csv}.}
\label{tab:m0-dic-ventas}
\begin{tabularx}{\textwidth}{@{}l l X@{}}
\toprule
\textbf{Columna} & \textbf{Tipo} & \textbf{Descripción} \\
\midrule
\codigo{fecha} & \codigo{<date>} & Día de la venta, del 2023-01-01 al
  2023-12-31 (uno por fila). \\
\codigo{dia\_semana} & \codigo{<chr>} & Nombre del día (``domingo'',
  ``lunes''\ldots). Depende del idioma del sistema al generar los datos. \\
\codigo{producto\_id} & \codigo{<dbl>} & Producto vendido, de 1 a 50.
  Llave foránea hacia \archivo{productos}. \\
\codigo{cliente\_id} & \codigo{<dbl>} & Cliente que compró, de 1 a 200.
  Llave foránea hacia \archivo{clientes}. \\
\codigo{cantidad} & \codigo{<dbl>} & Unidades vendidas, de 1 a 10. \\
\codigo{precio\_unitario} & \codigo{<dbl>} & Precio cobrado por unidad,
  aleatorio entre \$10 y \$500, con dos decimales. No coincide con el
  precio de catálogo del producto. \\
\codigo{descuento\_pct} & \codigo{<dbl>} & Descuento en porcentaje: 0, 5,
  10, 15 o 20 (con probabilidades de 50, 20, 15, 10 y 5\,\%). \\
\codigo{tienda\_id} & \codigo{<dbl>} & Tienda de la venta, de 1 a 10.
  Llave foránea hacia \archivo{tiendas}. \\
\codigo{vendedor\_id} & \codigo{<dbl>} & Empleado que vendió, de 1 a 25.
  Llave foránea hacia \codigo{empleado\_id} de \archivo{empleados}. \\
\codigo{subtotal} & \codigo{<dbl>} & \codigo{cantidad} $\times$
  \codigo{precio\_unitario}. \\
\codigo{descuento\_monto} & \codigo{<dbl>} & \codigo{subtotal} $\times$
  \codigo{descuento\_pct} / 100. \\
\codigo{total} & \codigo{<dbl>} & \codigo{subtotal} $-$
  \codigo{descuento\_monto}: lo que pagó el cliente. \\
\codigo{mes} & \codigo{<chr>} & Año y mes como texto: ``2023-01'' a
  ``2023-12''. \\
\codigo{trimestre} & \codigo{<chr>} & Trimestre: ``Q1'', ``Q2'', ``Q3'' o
  ``Q4''. \\
\bottomrule
\end{tabularx}
\end{table}

\subsection{\texorpdfstring{\archivo{clientes.csv}}{clientes.csv}: base de
clientes}

\textbf{200 filas y 10 columnas}. Cada fila es un cliente.

\begin{table}[H]
\centering\small
\caption{Diccionario de \archivo{clientes.csv}.}
\label{tab:m0-dic-clientes}
\begin{tabularx}{\textwidth}{@{}l l X@{}}
\toprule
\textbf{Columna} & \textbf{Tipo} & \textbf{Descripción} \\
\midrule
\codigo{cliente\_id} & \codigo{<dbl>} & Identificador único, de 1 a 200.
  Llave primaria. \\
\codigo{nombre} & \codigo{<chr>} & Nombre y apellido, combinados al azar
  de 20 nombres y 18 apellidos (puede repetirse). \\
\codigo{edad} & \codigo{<dbl>} & Edad en años, de 18 a 75. \\
\codigo{genero} & \codigo{<chr>} & ``M'' o ``F''. \\
\codigo{ciudad} & \codigo{<chr>} & Ciudad de residencia: CDMX,
  Guadalajara, Monterrey, Puebla, Querétaro, Tijuana, León, Mérida, Cancún
  o Toluca. \\
\codigo{fecha\_registro} & \codigo{<date>} & Fecha de alta como cliente,
  entre 2020-01-01 y 2023-12-31. \\
\codigo{es\_premium} & \codigo{<lgl>} & \codigo{TRUE} si pertenece al
  programa premium (alrededor de 30\,\% de los clientes). \\
\codigo{email} & \codigo{<chr>} & Correo en minúsculas con dominio
  \codigo{@email.com}. Se generó con otro nombre al azar, así que no
  corresponde a la columna \codigo{nombre}. \\
\codigo{telefono} & \codigo{<chr>} & Teléfono con formato
  ``55-XXXX-XXXX''. \\
\codigo{segmento} & \codigo{<chr>} & Segmento por edad: Joven (18 a 25),
  Adulto (26 a 40), Maduro (41 a 60) o Senior (61 a 75). \\
\bottomrule
\end{tabularx}
\end{table}

\subsection{\texorpdfstring{\archivo{productos.csv}}{productos.csv}:
catálogo de productos}

\textbf{50 filas y 13 columnas}. Cada fila es un producto del catálogo.

\begin{table}[H]
\centering\small
\caption{Diccionario de \archivo{productos.csv}.}
\label{tab:m0-dic-productos}
\begin{tabularx}{\textwidth}{@{}l l X@{}}
\toprule
\textbf{Columna} & \textbf{Tipo} & \textbf{Descripción} \\
\midrule
\codigo{producto\_id} & \codigo{<dbl>} & Identificador, de 1 a 50. Llave
  primaria. \\
\codigo{nombre\_producto} & \codigo{<chr>} & ``Producto 01'' a
  ``Producto 50''. \\
\codigo{categoria} & \codigo{<chr>} & Una de 10 categorías asignadas al
  azar: Electrónica, Ropa, Alimentos, Hogar, Deportes, Libros, Juguetes,
  Belleza, Automotriz o Mascotas. \\
\codigo{subcategoria} & \codigo{<chr>} & ``Sub 1'' a ``Sub 5''. \\
\codigo{marca} & \codigo{<chr>} & ``Marca A'' a ``Marca E''. \\
\codigo{precio\_catalogo} & \codigo{<dbl>} & Precio de lista, aleatorio
  entre \$15 y \$800. \\
\codigo{costo} & \codigo{<dbl>} & Costo unitario, aleatorio entre \$10 y
  \$400, generado sin relación con el precio. \\
\codigo{stock\_actual} & \codigo{<dbl>} & Unidades en inventario, de 0 a
  500. \\
\codigo{stock\_minimo} & \codigo{<dbl>} & Punto de reorden, de 10 a 50
  unidades. \\
\codigo{proveedor} & \codigo{<chr>} & ``Proveedor 1'' a ``Proveedor 4''. \\
\codigo{activo} & \codigo{<lgl>} & \codigo{TRUE} si el producto sigue a la
  venta (alrededor de 90\,\%). \\
\codigo{margen\_pct} & \codigo{<dbl>} & Margen bruto en porcentaje:
  (\codigo{precio\_catalogo} $-$ \codigo{costo}) /
  \codigo{precio\_catalogo} $\times$ 100, con dos decimales. Puede ser
  negativo. \\
\codigo{estado\_stock} & \codigo{<chr>} & ``Sin Stock'' si
  \codigo{stock\_actual} es 0; ``Stock Bajo'' si es menor que
  \codigo{stock\_minimo}; si no, ``Stock OK''. \\
\bottomrule
\end{tabularx}
\end{table}

\subsection{\texorpdfstring{\archivo{empleados.csv}}{empleados.csv}:
recursos humanos}

\textbf{100 filas y 12 columnas}. Cada fila es un empleado.

\begin{table}[H]
\centering\small
\caption{Diccionario de \archivo{empleados.csv}.}
\label{tab:m0-dic-empleados}
\begin{tabularx}{\textwidth}{@{}l l X@{}}
\toprule
\textbf{Columna} & \textbf{Tipo} & \textbf{Descripción} \\
\midrule
\codigo{empleado\_id} & \codigo{<dbl>} & Identificador, de 1 a 100. Llave
  primaria; el \codigo{vendedor\_id} de las ventas apunta aquí. \\
\codigo{nombre} & \codigo{<chr>} & Nombre y apellido al azar. \\
\codigo{departamento} & \codigo{<chr>} & Ventas, Marketing, IT, RRHH,
  Operaciones, Finanzas o Atención al Cliente. \\
\codigo{puesto} & \codigo{<chr>} & Analista, Especialista, Gerente,
  Coordinador o Asistente. \\
\codigo{salario\_mensual} & \codigo{<dbl>} & Salario mensual en pesos,
  aleatorio entre \$8,000 y \$45,000. \\
\codigo{fecha\_ingreso} & \codigo{<date>} & Fecha de contratación, entre
  2015-01-01 y 2023-12-31. \\
\codigo{sucursal} & \codigo{<chr>} & Matriz, Sucursal Norte, Sucursal Sur,
  Sucursal Este o Sucursal Oeste. \\
\codigo{genero} & \codigo{<chr>} & ``M'' o ``F''. \\
\codigo{edad} & \codigo{<dbl>} & Edad en años, de 22 a 65. \\
\codigo{nivel\_educacion} & \codigo{<chr>} & Preparatoria, Licenciatura,
  Maestría o Doctorado (con probabilidades de 20, 50, 25 y 5\,\%). \\
\codigo{antiguedad\_anos} & \codigo{<dbl>} & Años de antigüedad al
  2024-01-01 (días transcurridos / 365, con un decimal). \\
\codigo{salario\_anual} & \codigo{<dbl>} & \codigo{salario\_mensual}
  $\times$ 12. \\
\bottomrule
\end{tabularx}
\end{table}

\subsection{\texorpdfstring{\archivo{transacciones.csv}}{transacciones.csv}:
tickets de venta}

\textbf{1000 filas y 14 columnas}. Cada fila es un ticket de venta con su
hora y su método de pago. Es la tabla más parecida a la que exporta un
sistema de punto de venta real. Sus datos se generaron por separado de
\archivo{ventas\_retail.csv}, así que las dos tablas no ``cuadran'' entre
sí: úsalas como dos fuentes independientes.

\begin{table}[H]
\centering\small
\caption{Diccionario de \archivo{transacciones.csv}.}
\label{tab:m0-dic-transacciones}
\begin{tabularx}{\textwidth}{@{}l l X@{}}
\toprule
\textbf{Columna} & \textbf{Tipo} & \textbf{Descripción} \\
\midrule
\codigo{transaccion\_id} & \codigo{<dbl>} & Identificador del ticket, de 1
  a 1000. Llave primaria. \\
\codigo{fecha} & \codigo{<date>} & Día de la transacción, al azar dentro
  de 2023. \\
\codigo{hora} & \codigo{<time>} & Hora entre 8:00 y 20:45, en intervalos
  de 15 minutos. En el CSV está como texto sin cero inicial (``8:15'');
  \codigo{read\_csv()} la reconoce como hora. \\
\codigo{cliente\_id} & \codigo{<dbl>} & Cliente, de 1 a 200 (FK hacia
  \archivo{clientes}). \\
\codigo{producto\_id} & \codigo{<dbl>} & Producto, de 1 a 50 (FK hacia
  \archivo{productos}). \\
\codigo{cantidad} & \codigo{<dbl>} & Unidades, de 1 a 5. \\
\codigo{precio\_venta} & \codigo{<dbl>} & Precio unitario cobrado,
  aleatorio entre \$20 y \$600. \\
\codigo{metodo\_pago} & \codigo{<chr>} & Efectivo, Tarjeta Crédito,
  Tarjeta Débito o Transferencia (con probabilidades de 30, 35, 25 y
  10\,\%). \\
\codigo{tienda\_id} & \codigo{<dbl>} & Tienda, de 1 a 10 (FK hacia
  \archivo{tiendas}). \\
\codigo{vendedor\_id} & \codigo{<dbl>} & Vendedor, de 1 a 25 (FK hacia
  \codigo{empleado\_id}). \\
\codigo{total\_transaccion} & \codigo{<dbl>} & \codigo{cantidad} $\times$
  \codigo{precio\_venta} (esta tabla no tiene descuentos). \\
\codigo{mes} & \codigo{<chr>} & Año y mes: ``2023-01'' a ``2023-12''. \\
\codigo{trimestre} & \codigo{<chr>} & ``Q1'' a ``Q4''. \\
\codigo{dia\_semana} & \codigo{<chr>} & Nombre del día; depende del idioma
  del sistema. \\
\bottomrule
\end{tabularx}
\end{table}

\subsection{\texorpdfstring{\archivo{tiendas.csv}}{tiendas.csv}: red de
sucursales}

\textbf{10 filas y 8 columnas}. Cada fila es una tienda.

\begin{table}[H]
\centering\small
\caption{Diccionario de \archivo{tiendas.csv}.}
\label{tab:m0-dic-tiendas}
\begin{tabularx}{\textwidth}{@{}l l X@{}}
\toprule
\textbf{Columna} & \textbf{Tipo} & \textbf{Descripción} \\
\midrule
\codigo{tienda\_id} & \codigo{<dbl>} & Identificador, de 1 a 10. Llave
  primaria. \\
\codigo{nombre\_tienda} & \codigo{<chr>} & ``Sucursal Centro'', ``Sucursal
  Norte'', \ldots, ``Sucursal Express''. \\
\codigo{ciudad} & \codigo{<chr>} & CDMX (4 tiendas), Guadalajara (2),
  Monterrey, Puebla, Tijuana y Querétaro (1 cada una). \\
\codigo{zona} & \codigo{<chr>} & Centro, Norte, Sur, Este u Oeste. \\
\codigo{tamano\_m2} & \codigo{<dbl>} & Superficie de la tienda en metros
  cuadrados, de 300 a 800. \\
\codigo{empleados} & \codigo{<dbl>} & Número de empleados de la tienda,
  de 12 a 40. No es una llave hacia \archivo{empleados.csv}. \\
\codigo{fecha\_apertura} & \codigo{<date>} & Fecha de apertura, entre
  2009-01-15 y 2020-07-01. \\
\codigo{tipo} & \codigo{<chr>} & Principal, Estándar, Outlet, Premium o
  Express. \\
\bottomrule
\end{tabularx}
\end{table}

\subsection{\texorpdfstring{\archivo{datos\_empresa.xlsx}}{datos\_empresa.xlsx}:
el mismo negocio en Excel}

Un libro de Excel con cinco hojas, para practicar la importación de
archivos de Excel en el Módulo~\ref{cap:fundamentos}: \textbf{Ventas} (las
primeras 100 filas de \archivo{ventas\_retail}), \textbf{Clientes} (los
primeros 50), \textbf{Productos} (los 50), \textbf{Empleados} (los primeros
50) y \textbf{Tiendas} (las 10).

\subsection{Datos simulados: lo que debes saber}
\label{sec:m0-calidad}

Los datos son \textbf{simulados}: se generaron al azar con reglas simples,
no provienen de una empresa real. Eso tiene una consecuencia que no vamos a
esconder: algunas columnas se generaron de forma independiente y producen
combinaciones que en un negocio real serían errores. Lejos de ser un
problema, es una oportunidad, porque los datos reales tampoco vienen
limpios. Revisar la \termino{calidad de datos} antes de analizar es una de
las tareas más importantes de un analista, y aquí la practicarás desde el
primer día (tabla~\ref{tab:m0-calidad}).

\begin{table}[H]
\centering\small
\caption{Hallazgos de calidad de datos conocidos en los datos del curso.}
\label{tab:m0-calidad}
\begin{tabularx}{\textwidth}{@{}>{\raggedright\arraybackslash}X
    >{\raggedright\arraybackslash}X@{}}
\toprule
\textbf{Hallazgo} & \textbf{Por qué ocurre} \\
\midrule
Productos con margen negativo (costo mayor que precio). &
  \codigo{precio\_catalogo} y \codigo{costo} se generaron por separado. \\
El precio cobrado en las ventas no coincide con el precio de catálogo. &
  \codigo{precio\_unitario} y \codigo{precio\_venta} son aleatorios. \\
Ventas de productos marcados como inactivos. &
  \codigo{activo} se asignó al azar, sin mirar las ventas. \\
Ventas anteriores a la fecha de registro del cliente. &
  \codigo{fecha\_registro} llega hasta diciembre de 2023. \\
El correo de un cliente no corresponde a su nombre. &
  \codigo{email} se armó con otro nombre al azar. \\
Hay ``vendedores'' que no son del departamento de Ventas. &
  \codigo{vendedor\_id} toma los empleados 1 a 25, sin importar su
  departamento. \\
Una sola venta por día en \archivo{ventas\_retail}, y no cuadra con
  \archivo{transacciones}. &
  Son dos muestras independientes y simplificadas. \\
La suma de \codigo{empleados} en \archivo{tiendas} no coincide con las 100
  filas de \archivo{empleados.csv}, y la sucursal ``Matriz'' no existe en
  \archivo{tiendas}. &
  Las dos tablas se generaron sin relación entre sí. \\
Hay clientes en ciudades sin tienda (León, Mérida, Cancún, Toluca). &
  Posible en la realidad (compran de visita), pero conviene saberlo. \\
\bottomrule
\end{tabularx}
\end{table}

No necesitas creernos: el siguiente código comprueba varios de estos
hallazgos. Usa funciones que aprenderás en el
Módulo~\ref{cap:manipulacion}, así que por ahora basta con que leas los
comentarios. Reutiliza la tabla \codigo{ventas} que leíste en tu primer
análisis.

\begin{rcode}
# Leemos las tablas que vamos a revisar
productos <- read_csv("datasets/productos.csv", show_col_types = FALSE)
clientes  <- read_csv("datasets/clientes.csv", show_col_types = FALSE)
empleados <- read_csv("datasets/empleados.csv", show_col_types = FALSE)

# Hallazgo 1: ¿cuántos productos tienen margen negativo?
sum(productos$margen_pct < 0)

# Hallazgo 2: ¿cuántas ventas son de productos inactivos?
ventas %>%
  left_join(productos, by = "producto_id") %>%  # agrega datos del producto
  filter(activo == FALSE) %>%                   # solo productos inactivos
  nrow()                                        # cuenta las filas

# Hallazgo 3: ¿cuántas ventas ocurrieron ANTES del registro del cliente?
ventas %>%
  left_join(clientes, by = "cliente_id") %>%    # agrega datos del cliente
  filter(fecha < fecha_registro) %>%            # compra antes del alta
  nrow()
\end{rcode}

\begin{rsalida}
[1] 10
[1] 65
[1] 42
\end{rsalida}

De los 50 productos, 10 tienen margen negativo; 65 de las 365 ventas
corresponden a productos inactivos, y en 42 ventas el cliente ``compró''
antes de registrarse. Ahora los dos hallazgos que se ven a simple vista:

\begin{rcode}
# Hallazgo 4: el nombre y el correo no corresponden a la misma persona
clientes %>%
  select(cliente_id, nombre, email) %>%
  head(3)

# Hallazgo 5: ¿de qué departamento son los "vendedores" 1 a 25?
empleados %>%
  filter(empleado_id <= 25) %>%
  count(departamento, sort = TRUE)
\end{rcode}

\begin{rsalida}
# A tibble: 3 x 3
  cliente_id nombre          email
       <dbl> <chr>           <chr>
1          1 Ricardo Morales juan.ramírez@email.com
2          2 Miguel Cruz     alberto.jiménez@email.com
3          3 Miguel Flores   pedro.sánchez@email.com
# A tibble: 7 x 2
  departamento            n
  <chr>               <int>
1 Atención al Cliente    10
2 RRHH                    4
3 Ventas                  4
4 Marketing               3
5 Finanzas                2
6 IT                      1
7 Operaciones             1
\end{rsalida}

El cliente 1 se llama Ricardo Morales, pero su correo es de ``juan.ramírez''.
Y de los 25 empleados que aparecen como vendedores, solo 4 pertenecen al
departamento de Ventas. En una empresa real, cada uno de estos hallazgos
sería una conversación con el dueño del sistema de origen. En el curso los
usaremos para practicar: los detectarás, decidirás cómo tratarlos y
documentarás tu decisión.

\begin{advertencia}
Por lo mismo, las cifras del curso no representan a ninguna empresa real.
No saques conclusiones de negocio sobre el comercio en México a partir de
ellas: son un campo de práctica. Lo que sí es real es el método, que
aplicarás igual con los datos de tu trabajo.
\end{advertencia}

\begin{hazlotu}
En el panel \textit{Files} de RStudio haz clic en
\archivo{datasets/productos.csv} y elige \menu{View File}. Busca el
producto con el margen más negativo. ¿Qué pasaría con un reporte de
rentabilidad por categoría si nadie lo detectara? Pista: en la salida de
arriba, R te dijo cuántos productos revisar.
\end{hazlotu}

% ----------------------------------------------------------------------------
\section{El código del libro y cómo compilarlo}
\label{sec:m0-codigo}
% ----------------------------------------------------------------------------

\subsection{Los scripts de la carpeta \texorpdfstring{\archivo{codigo/}}{codigo/}}

La carpeta \archivo{latex/codigo/} contiene un script de R por módulo (por
ejemplo, \archivo{01-fundamentos.R} para el Módulo~\ref{cap:fundamentos})
con todo el código ejecutable de ese módulo, en el mismo orden que en el libro.
Los bloques de franja naranja (instalaciones, apps Shiny, conexiones a
servidores) aparecen comentados, solo como referencia. Para usarlos:

\begin{enumerate}
  \item Abre tu proyecto \archivo{curso-bi-rstudio.Rproj}, para que el
        directorio de trabajo sea la carpeta del curso.
  \item Abre el script del módulo desde el panel \textit{Files}.
  \item Ejecútalo bloque por bloque con \tecla{Ctrl} + \tecla{Enter} mientras
        lees el módulo. Cada bloque está marcado con un comentario
        \codigo{\# ---- Bloque N}.
\end{enumerate}

Estos scripts se generan automáticamente a partir de los capítulos del
libro con el programa \archivo{herramientas/extraer\_codigo.py} (escrito en
Python), así que siempre coinciden con lo impreso. Si modificas un capítulo
y quieres regenerarlos, desde la carpeta \archivo{latex/} ejecuta en una
terminal:

\begin{textocodigo}
python3 herramientas/extraer_codigo.py
\end{textocodigo}

\subsection{Compilar el libro}

El libro está escrito en \LaTeX{}. Si quieres generar tu propio PDF (por
ejemplo, después de añadir notas a un capítulo), tienes dos caminos:

\begin{itemize}
  \item \textbf{Overleaf} (en línea, sin instalar nada): comprime la
        carpeta \archivo{latex/} en un ZIP, crea un proyecto con
        \menu{New Project > Upload Project}, verifica en \menu{Menu} que el
        compilador sea \textbf{pdfLaTeX} y que el documento principal sea
        \archivo{curso-bi-rstudio.tex}, y pulsa \menu{Recompile}. Si el
        plan gratuito se queda sin tiempo de compilación, compila en tu
        computadora.
  \item \textbf{En tu computadora}, con una distribución de \LaTeX{} como
        TeX Live (Windows y Linux) o MacTeX (macOS), abre una terminal en
        la carpeta \archivo{latex/} y ejecuta:
\end{itemize}

\begin{textocodigo}
latexmk -pdf curso-bi-rstudio.tex
\end{textocodigo}

\codigo{latexmk} compila las veces que haga falta para que el índice y las
referencias cruzadas (``Solución en la página\ldots'') queden correctos.

% ----------------------------------------------------------------------------
\section{Consejos para aprender}
\label{sec:m0-consejos}
% ----------------------------------------------------------------------------

Aprender R es como aprender un idioma: se aprende hablándolo. Estos
consejos vienen de muchas generaciones de estudiantes:

\begin{enumerate}
  \item \textbf{Escribe el código tú mismo}, no lo copies. Teclear cada
        línea te obliga a leerla, y los errores que cometas al escribir te
        enseñarán a leer los mensajes de R.
  \item \textbf{Predice antes de ejecutar.} Antes de pulsar \tecla{Ctrl} +
        \tecla{Enter}, pregúntate qué va a salir. Si aciertas, lo
        entendiste; si no, acabas de aprender algo.
  \item \textbf{Lee los errores con calma.} El mensaje casi siempre dice qué
        pasó y en qué función. Los más comunes son un paréntesis o una
        comilla sin cerrar, una coma de más, un nombre mal escrito o un
        paquete sin cargar.
  \item \textbf{Rompe el código a propósito.} Cambia un número, quita un
        argumento, usa otra columna. Ver qué cambia es la forma más rápida
        de entender qué hace cada parte.
  \item \textbf{Practica con tus propios datos.} Cuando termines un módulo,
        repite un ejemplo con datos de tu trabajo (exportados a CSV y
        respetando siempre la confidencialidad). Nada motiva más que
        resolver un problema real.
  \item \textbf{Constancia antes que intensidad.} Una hora diaria rinde
        mucho más que siete horas el domingo.
  \item \textbf{Intenta los ejercicios antes de ver la solución.} Dale al
        menos 20 minutos a cada uno. Si te atoras, lee solo la primera
        línea de la solución y vuelve a intentarlo.
  \item \textbf{Lleva tu propio acordeón}: un script con las funciones que
        vas aprendiendo y un ejemplo corto de cada una. El
        Apéndice~\ref{ap:referencia} te da una base.
  \item \textbf{Explica lo que aprendiste} a un colega. Si puedes
        explicarlo con palabras sencillas, lo dominas.
\end{enumerate}

% ----------------------------------------------------------------------------
\section{Cómo pedir ayuda}
\label{sec:m0-ayuda}
% ----------------------------------------------------------------------------

\subsection{La ayuda dentro de R}

Toda función de R tiene una página de ayuda con su descripción, sus
argumentos y ejemplos al final. Es la primera parada cuando no recuerdas
cómo se usa algo:

\begin{rnoejecutar}
?mean                      # ayuda de una función (o F1 sobre su nombre)
help("read_csv")           # lo mismo, escrito de otra forma
??"regression"             # busca un tema en todas las ayudas instaladas
example(mean)              # ejecuta los ejemplos de la ayuda
args(round)                # muestra solo los argumentos de una función
vignette("dplyr")          # guía larga (viñeta) de un paquete
\end{rnoejecutar}

La ayuda de R está en inglés. Léela primero por la sección
\textit{Examples}, al final: suele ser lo más claro.

\begin{consejo}
Cuando busques un error en internet, búscalo en inglés: hay muchas más
respuestas. Si tu R muestra los mensajes en español, puedes cambiarlos a
inglés en la sesión actual con \codigo{Sys.setenv(LANGUAGE = "en")}, volver
a ejecutar la línea que falló y copiar el mensaje exacto en el buscador.
\end{consejo}

\subsection{Recursos en línea}

\begin{itemize}
  \item \textbf{Hojas de referencia (\textit{cheatsheets}) de Posit}: una o
        dos páginas por paquete con las funciones principales; varias están
        traducidas al español. Desde RStudio: \menu{Help > Cheat Sheets}, o
        en \url{https://posit.co/resources/cheatsheets/}.
  \item \textbf{R for Data Science}, de Hadley Wickham y colaboradores,
        gratuito en línea (\url{https://r4ds.hadley.nz}). Su primera
        edición está traducida al español como \textit{R para Ciencia de
        Datos} (\url{https://es.r4ds.hadley.nz}).
  \item \textbf{Stack Overflow}, el sitio de preguntas y respuestas de
        programación: en inglés con la etiqueta \codigo{r}
        (\url{https://stackoverflow.com/questions/tagged/r}) y en español
        en \url{https://es.stackoverflow.com}.
  \item \textbf{Posit Community}, el foro oficial de RStudio y el
        tidyverse: \url{https://forum.posit.co}.
  \item \textbf{TidyTuesday}: cada semana la comunidad publica un conjunto
        de datos nuevo para practicar y compartir análisis
        (\url{https://github.com/rfordatascience/tidytuesday}).
  \item \textbf{R-bloggers}: artículos y tutoriales de toda la comunidad
        (\url{https://www.r-bloggers.com}).
\end{itemize}

\subsection{Comunidad de R en español}

R tiene una comunidad hispanohablante muy activa y acogedora con los
principiantes:

\begin{itemize}
  \item \textbf{R-Ladies} (\url{https://rladies.org}): organización mundial
        que promueve la diversidad de género en la comunidad de R, con
        capítulos locales en muchas ciudades de México y Latinoamérica.
        Organiza reuniones y talleres gratuitos, abiertos a todas las
        personas interesadas.
  \item \textbf{Comunidad R Hispano} (\url{https://comunidadrhispano.es}):
        asociación de usuarios de R en español; organiza las Jornadas de R
        en España.
  \item \textbf{LatinR} (\url{https://latin-r.com}): conferencia
        latinoamericana sobre uso de R en investigación y desarrollo, con
        charlas y talleres en español y portugués.
  \item \textbf{Grupos locales de usuarios de R}: busca en Meetup o en
        redes sociales ``R user group'' o ``useR'' seguido de tu ciudad.
\end{itemize}

\subsection{Cómo hacer una buena pregunta}

Quien te ayuda no ve tu pantalla. Una buena pregunta incluye un
\termino{ejemplo reproducible} (\textit{reprex}): el código mínimo que
produce el problema, unos pocos datos de ejemplo (creados en el propio
código, nunca datos confidenciales de tu empresa), el mensaje de error
completo y lo que esperabas obtener. Muchas veces, al preparar el ejemplo
mínimo, encontrarás tú mismo la solución. El paquete \paquete{reprex}, que
se instala con el tidyverse, te ayuda a preparar ese ejemplo con formato
listo para pegar en un foro.

Los asistentes de inteligencia artificial también pueden explicarte un
mensaje de error o sugerirte código. Úsalos como a un compañero que a veces
se equivoca: ejecuta lo que te propongan, comprueba que el resultado tenga
sentido y asegúrate de entender cada línea antes de usarla en un reporte.

\begin{resumen}[title={Resumen de la introducción}]
\begin{itemize}
  \item \textbf{BI} convierte datos en información para decidir;
        \textbf{Business Analytics} añade explicación, predicción y
        recomendación. Hay cuatro tipos de analítica: descriptiva,
        diagnóstica, predictiva y prescriptiva.
  \item R y RStudio aportan análisis \textbf{reproducible},
        estadística, gráficas de calidad, dashboards y reportes
        automáticos. Complementan a Excel, Power BI y Python.
  \item Todo proyecto de BI sigue un ciclo: pregunta, datos, ETL, análisis,
        comunicación, decisión y medición.
  \item Los bloques grises (\codigo{rcode}) se ejecutan; los de franja
        naranja los ejecutas tú cuando haga falta; las salidas tienen
        franja azul verdosa. El \codigo{[1]} es un índice, no un dato.
  \end{itemize}
\textbf{Lista de verificación antes del Módulo~\ref{cap:fundamentos}:}
\begin{itemize}
  \item[\ding{113}] R y RStudio instalados y configurados (sin guardar
        \archivo{.RData}).
  \item[\ding{113}] Paquetes del curso instalados.
  \item[\ding{113}] Proyecto \archivo{curso-bi-rstudio.Rproj} creado.
  \item[\ding{113}] Datos generados con \codigo{source("datasets/generar\_datasets.R")}.
  \item[\ding{113}] Script de verificación sin líneas \codigo{FALTA}.
  \item[\ding{113}] Tu primer análisis (ventas por trimestre) ejecutado.
\end{itemize}
\end{resumen}

¿Todo listo? Entonces abre tu proyecto y pasa al
Módulo~\ref{cap:fundamentos}. Empezamos por lo más básico: cómo piensa R.

\endgroup
